\documentclass[11pt, a4paper]{article}

% ── Packages ──────────────────────────────────────────────────────────────────
\usepackage[utf8]{inputenc}
\usepackage[T1]{fontenc}

\usepackage{amsmath, amssymb, amsthm}
\usepackage{mathtools}
\usepackage{geometry}
\usepackage{setspace}
\usepackage{booktabs}
\usepackage{array}
\usepackage{longtable}
\usepackage{hyperref}
\usepackage{cleveref}
\usepackage{enumitem}
\usepackage{titlesec}
\usepackage{abstract}
\usepackage{fancyhdr}
\usepackage{xcolor}
\usepackage{listings}
\usepackage{tikz-cd}

% ── Page geometry ─────────────────────────────────────────────────────────────
\geometry{
  top    = 2.8cm,
  bottom = 2.8cm,
  left   = 3.0cm,
  right  = 3.0cm,
  headheight = 14pt
}

% ── Spacing ───────────────────────────────────────────────────────────────────
\onehalfspacing

% ── Header / footer ───────────────────────────────────────────────────────────
\pagestyle{fancy}
\fancyhf{}
\rhead{\small RSVP and the MOND Transition}
\lhead{\small Flyxion}
\cfoot{\thepage}
\renewcommand{\headrulewidth}{0.4pt}

% ── Listings style for code/grammar ──────────────────────────────────────────
\lstset{
  basicstyle=\small\ttfamily,
  breaklines=true,
  frame=single,
  xleftmargin=1em,
  xrightmargin=1em,
  numbers=none,
  escapeinside={(*}{*)}
}

% ── Theorem environments ──────────────────────────────────────────────────────
\theoremstyle{plain}
\newtheorem{theorem}{Theorem}[section]
\newtheorem{conjecture}[theorem]{Conjecture}
\newtheorem{corollary}[theorem]{Corollary}
\newtheorem{lemma}[theorem]{Lemma}

\theoremstyle{definition}
\newtheorem{definition}[theorem]{Definition}
\newtheorem{tschema}[theorem]{Theorem-Schema}
\newtheorem{axiom}[theorem]{Axiom}
\newtheorem{heuristic}[theorem]{Heuristic}

\theoremstyle{remark}
\newtheorem{remark}[theorem]{Remark}
\newtheorem{obligation}[theorem]{Open Proof Obligation}
\newtheorem{observation}[theorem]{Observation}

% ── Hyperref ──────────────────────────────────────────────────────────────────
\hypersetup{
  colorlinks = true,
  linkcolor  = blue!60!black,
  citecolor  = blue!60!black,
  urlcolor   = blue!60!black,
  pdftitle   = {Admissibility Geometry and the MOND Transition},
  pdfauthor  = {Flyxion}
}

% ── Custom commands ───────────────────────────────────────────────────────────
\newcommand{\Xfield}{X = (\Phi,\,\mathbf{v},\,S)}
\newcommand{\ctrl}{\tfrac{S}{\Phi}}
\newcommand{\ctrlell}{\chi_\ell}
\newcommand{\RA}{R_{\!A}}
\newcommand{\aA}{a_{\!A}}
\newcommand{\sheaf}{\mathcal{F}}
\newcommand{\cover}{\{U_\alpha\}}
\newcommand{\recon}{\mathcal{R}}
\newcommand{\Ucover}{\mathcal{U}}
\newcommand{\Rcoarse}{\mathcal{R}_\lambda}
\newcommand{\Pitherm}{\Pi_{\mathrm{therm}}}
\newcommand{\Piinfo}{\Pi_{\mathrm{info}}}
\newcommand{\Pigrav}{\Pi_{\mathrm{grav}}}
\newcommand{\Piconf}{\Pi_{\mathrm{conf}}}
\newcommand{\Hcech}{\check{H}}
\newcommand{\Ccech}{\check{C}}

% ── Title ─────────────────────────────────────────────────────────────────────
\title{%
  \textbf{Admissibility Geometry and the MOND Transition}\\[0.4em]
  \large Reconstruction Obstruction as an Alternative to Dark Matter
}
\author{Flyxion}
\date{\today}

% ══════════════════════════════════════════════════════════════════════════════
\begin{document}
% ══════════════════════════════════════════════════════════════════════════════

\maketitle
\thispagestyle{empty}

% ── Abstract ──────────────────────────────────────────────────────────────────
\begin{abstract}
\noindent
We develop a reinterpretation of gravitational anomalies --- including flat rotation
curves, weak lensing excesses, and the MOND acceleration scale --- within the
Relativistic Scalar-Vector Plenum (RSVP) framework. Existing emergent gravity
programs derive spacetime geometry from physical exchange processes, thermodynamic
bookkeeping, or information flow. RSVP inverts this order: admissibility geometry
is treated as primary, and persistent physical exchange is interpreted as a
stabilized reconstruction mode within that geometry.

The framework introduces the coupled field $\Xfield$, where $\Phi$ encodes the
accessibility geometry of continuation space, $\mathbf{v}$ encodes directed
coherence flow, and $S$ measures admissibility dispersion. The dimensionless ratio
$S/\Phi$ serves as a central control parameter governing local reconstruction
factorization. We define an \emph{admissibility radius} $\RA$ as the largest scale
over which local reconstruction remains approximately factorizable without
cosmological boundary data, and conjecture that the RSVP field equations determine
$\RA \propto c/H_0$ in weak-gradient, low-density regimes, yielding the MOND
acceleration scale $a_0 \sim cH_0$ as a geometric consequence of reconstruction
structure rather than a modification of force law.

The paper operates across four explicitly distinguished epistemic registers:
definitions, structural conjectures, derived implications, and open proof
obligations. We do not derive MOND from completed RSVP field equations; we identify
the geometric objects from which such a derivation would have to proceed. A dedicated
section identifies the five irreducible commitments of the framework, separating
load-bearing assumptions from illustrative scaffolding. Appendices
develop the formal sheaf structure, constraint propagation, coarse-graining flow with
attractor analysis, candidate dynamical equations, obstruction cohomology, relation
to renormalization theory, observational targets, locality clarification, formal
implementation via BNF grammar and Lean-style schemas, a variational formulation
sketch, effective metric emergence, and admissibility gauge equivalence.
\end{abstract}

\vspace{1em}
\noindent\textit{Keywords:} emergent gravity, MOND, admissibility geometry,
reconstruction sheaf, entropic gravity, dark matter, RSVP framework, obstruction
cohomology, renormalization flow.

\newpage

% ── Notation table ────────────────────────────────────────────────────────────
\subsection*{Core Objects and Notation}

\begin{center}
\renewcommand{\arraystretch}{1.4}
\begin{tabular}{@{}lp{9cm}@{}}
\toprule
\textbf{Symbol} & \textbf{Meaning} \\
\midrule
$X = (\Phi, \mathbf{v}, S)$ & Coupled RSVP field \\
$\Phi$ & Continuation accessibility scalar field \\
$\mathbf{v}$ & Coherence-flow vector field \\
$S$ & Admissibility dispersion (entropy) field \\
$\chi = S/\Phi$ & Reconstruction control parameter \\
$\chi_\ell = S_\ell/\Phi_\ell$ & Scale-dependent control parameter \\
$\Ucover = \cover$ & Reconstruction (admissibility) cover \\
$\sheaf$ & Reconstruction sheaf \\
$\sheaf(U_\alpha)$ & Admissible sections over neighborhood $U_\alpha$ \\
$\rho_{\alpha\beta}$ & Restriction map $\sheaf(U_\alpha) \to \sheaf(U_\alpha \cap U_\beta)$ \\
$\RA$ & Admissibility radius \\
$\aA$ & Admissibility acceleration scale $\sim c^2/\RA$ \\
$\mathcal{P}[s]$ & Persistence functional of section $s$ \\
$\Rcoarse$ & Coarse-graining operator on covers \\
$\Omega_\ell[s]$ & Scale-dependent obstruction functional \\
$\beta(\chi)$ & Admissibility beta-function \\
$\Pitherm,\Piinfo,\Pigrav,\Piconf$ & Entropy projection functors \\
$\Hcech^k(\Ucover, \sheaf)$ & \v{C}ech cohomology of cover with sheaf coefficients \\
\bottomrule
\end{tabular}
\end{center}

\newpage
\tableofcontents
\newpage

% ══════════════════════════════════════════════════════════════════════════════
\section{Introduction}
% ══════════════════════════════════════════════════════════════════════════════

The relationship between gravity and thermodynamics has attracted sustained
theoretical attention since Bekenstein and Hawking established that black holes carry
entropy proportional to horizon area. Subsequent programs --- Jacobson's derivation
of the Einstein equations from thermodynamic identities \cite{jacobson1995,jacobson2016},
Verlinde's entropic gravity \cite{verlinde2011,verlinde2017}, and more recently
transactional entropic gravity \cite{kastner2024} --- have attempted to establish
that gravitational dynamics emerge from statistical or information-theoretic structure
rather than constituting a fundamental interaction. Padmanabhan has developed related
thermodynamic perspectives on gravity \cite{padmanabhan2010}, and the connection
between entanglement structure and spacetime geometry has been explored through
holographic frameworks \cite{vanraamsdonk2010,swingle2012}.

These programs share a common derivation order: physical exchange processes,
thermodynamic bookkeeping, or entropic flow are treated as primitive, and spacetime
geometry is reconstructed from them. The present work inverts this order. We treat
\emph{admissibility geometry} as primary and interpret persistent physical exchange
as a stabilized reconstruction mode within that geometry. Emitters, absorbers,
particles, and force laws are not inputs to the framework but outputs --- roles
stabilized by recursive cross-scale coherence of the coupled field $\Xfield$ within
its reconstruction sheaf \cite{bredon1997,maclane1992}.

This inversion has a specific consequence for the MOND anomaly
\cite{milgrom1983,famaey2012}. Existing programs explain the transition scale
$a_0 \sim cH_0$ by finding thermodynamic narratives that reproduce it. The present
framework instead asks: what geometric property of the reconstruction sheaf would
make such a scale \emph{structurally necessary}? The answer we propose --- an
admissibility radius $\RA$ beyond which local reconstruction requires cosmological
gluing data --- converts a phenomenological constant into a constraint on the field
theory itself.

\subsection*{Epistemic Registers}

The paper operates simultaneously across four epistemic registers which must be
carefully distinguished. Certain structures are \textbf{definitional}, specifying
the reconstruction geometry and admissibility framework. Others are
\textbf{conjectural}, proposing relationships between admissibility structure and
cosmological scaling. Still others are \textbf{derived implications} under specified
assumptions. Finally, several results remain \textbf{open proof obligations}
generated by the framework itself. We make these distinctions syntactically explicit
throughout --- using the theorem environments \emph{Definition}, \emph{Conjecture},
\emph{Theorem-Schema}, \emph{Heuristic}, and \emph{Open Proof Obligation} --- in
order to avoid conflating motivation, formal consequence, and unresolved structure, a
conflation that has weakened much of the emergent gravity literature.

\subsection*{Structure of the Paper}

Section~\ref{sec:geometry} develops the admissibility geometry and sheaf structure,
including an explicit clarification of the operational meaning of the RSVP fields.
Section~\ref{sec:transition} introduces the admissibility phase transition, scaling
constraints, and the central conjecture regarding $\RA$.
Section~\ref{sec:entropy} treats entropy species as regime projections including
commutator structure.  Section~\ref{sec:darkmatter} reinterprets dark matter
phenomenology as reconstruction obstruction.  Section~\ref{sec:comparison} compares
the framework with neighboring programs.  Section~\ref{sec:assumptions} identifies
the five irreducible commitments of the theory.
Sections~\ref{sec:inversion}--\ref{sec:completion} develop the philosophical
inversion, regime interpretation, observational proxies, forbidden continuation
classes, structural risks, and completion criteria.
Section~\ref{sec:conclusion} consolidates contributions and open
obligations. Appendices~\ref{app:sheaf}--\ref{app:gauge} develop the formal
machinery: sheaf structure with constraint propagation, coarse-graining flow with
attractor analysis and coarse-graining disambiguation, candidate field equations,
obstruction cohomology, renormalization analogy, observational targets, locality
clarification, formal implementation, a variational formulation sketch, effective
metric emergence, and admissibility gauge equivalence.

% ══════════════════════════════════════════════════════════════════════════════
\section{Admissibility Geometry and the Reconstruction Sheaf}
\label{sec:geometry}
% ══════════════════════════════════════════════════════════════════════════════

\subsection{The Coupled Field}

The fundamental object of the RSVP framework is the coupled field
\begin{equation}
  \Xfield,
\end{equation}
where $\Phi$ is a scalar field encoding the \emph{accessibility geometry of
continuation space} --- the structure governing which local field trajectories admit
coherent extension into persistent reconstruction modes; $\mathbf{v}$ is a vector
field encoding directed transport and coherence flow; and $S$ is a scalar
entropy-dispersion field measuring the spread of admissible reconstruction
trajectories.

These three components are dynamically coupled. The scalar field $\Phi$ governs
which continuations remain accessible; $\mathbf{v}$ tracks preferred directions of
coherent continuation; and $S$ measures how tightly reconstruction trajectories
concentrate near stable modes. A configuration with high $S$ relative to $\Phi$ has
many admissible continuations and correspondingly weak reconstruction stability. Low
$S/\Phi$ concentrates trajectories near persistent modes --- the regime in which
particles, clocks, and force-mediating exchanges stabilize as recognizable roles.

\begin{remark}[Operational meaning of the RSVP fields]
\label{rem:operational}
The field $\Phi$ is \emph{not} identified with the Newtonian gravitational potential,
nor with any standard matter field. It encodes a more primitive object: the geometry
of which local continuations remain stably extendable. The Newtonian potential
emerges as a low-$\chi$ effective reconstruction observable, in the same way that
hydrodynamic pressure is not microscopically fundamental but emerges from kinetic
theory in appropriate regimes. Similarly, $S$ should not be read primarily as
thermodynamic entropy; it is admissibility dispersion, and thermodynamic entropy is
only one of its projection functors ($\Pitherm$) under specific coarse-grainings.
Foregrounding the term \emph{dispersion} rather than \emph{entropy} for $S$ would
reduce confusion with thermodynamic baggage; the present paper retains both usages
but always interprets $S$ through the projection-functor framework rather than as a
primitive thermodynamic quantity.
\end{remark}

\medskip
\noindent\textbf{The control parameter.}\quad The dimensionless ratio
\begin{equation}
  \label{eq:control}
  \chi \;=\; \ctrl
\end{equation}
is the central control parameter of the framework. When $\chi \ll 1$,
entropy-dispersion is small relative to continuation accessibility, local
reconstruction is strongly factorizable, and effective dynamics reduce approximately
to Newtonian behavior. When $\chi \sim 1$, dispersion becomes comparable to
accessibility, local factorization breaks down, and reconstruction requires boundary
data from larger scales to maintain global stability. This ratio governs
simultaneously: local factorization stability, the admissibility phase transition,
entropy-channel convergence, and obstruction growth dynamics. Its ubiquity across all
four phenomena is evidence of the framework's internal coherence rather than
coincidence.

\subsection{Admissible Local Sections}

\begin{definition}[Admissible local section]
\label{def:admissible}
A local section of $X$ over a region $U$ is \emph{admissible} if reconstruction
trajectories originating within $U$ remain dynamically stable under bounded
perturbation and admit coherent continuation under the restriction structure of the
field cover.
\end{definition}

Admissibility is a \emph{persistence condition}: it asks not merely whether a
configuration is locally self-consistent, but whether it can continue coherently into
neighboring regions and across scales. Three notions must be carefully distinguished.

\begin{enumerate}[label=(\roman*)]
  \item \textbf{Local dynamical stability.} Trajectories within $U$ remain bounded
    under perturbation. This is a condition on the field equations restricted to $U$
    alone, making no reference to neighboring regions.

  \item \textbf{Global admissibility.} Locally stable sections admit coherent
    extension across the full reconstruction cover without generating contradictions
    at overlapping boundaries. A section may be locally stable while failing global
    admissibility if its continuation structure conflicts with that of neighboring
    sections.

  \item \textbf{Reconstruction persistence.} Globally admissible extensions remain
    stable under renormalization of the reconstruction cover across scales --- a
    multiscale stability condition rather than merely a temporal one. A section may
    be globally admissible at one scale while failing persistence at another if
    $S$-dispersion grows faster than the coherence structure of $\mathbf{v}$ can
    absorb under coarse-graining.
\end{enumerate}

The distinction between local stability and global admissibility is precisely what
separates RSVP from transactional approaches \cite{kastner2024}. Transactional
completion guarantees local interaction closure --- offer and confirmation waves match
at a spacetime event. RSVP admissibility requires persistent global extension
compatibility across the entire reconstruction cover. A transaction can complete
locally while generating globally obstructed reconstruction dynamics; RSVP classifies
such processes as inadmissible or metastable regardless of their local consistency.

\subsection{The Reconstruction Cover}

\begin{definition}[Admissibility neighborhood]
An \emph{admissibility neighborhood} $U_\alpha$ is a region over which
reconstruction remains approximately factorizable, meaning the admissible section
over $U_\alpha$ can be specified without reference to data outside $U_\alpha$ up to
corrections of order $\ctrl$.
\end{definition}

We introduce a cover $\Ucover = \cover$ of the field domain where each $U_\alpha$
is an admissibility neighborhood. Regions where $S/\Phi$ is small generate large,
strongly factorizable neighborhoods. Regions where $S/\Phi$ approaches unity
generate small, weakly factorizable neighborhoods whose sections increasingly require
external boundary data for stable continuation.

\begin{remark}
\label{rem:cover-derived}
The reconstruction cover is a \emph{derived} object of the field geometry, not a
background against which the field is defined. Classical spacetime structure, to the
extent it appears, is a stable regime of the reconstruction cover rather than its
presupposition. This is the formal content of RSVP's departure from ordinary field
theory: integrating out short-scale reconstruction structure changes not merely
effective couplings but the factorization structure of admissible reconstruction
itself \cite{wilson1974,goldenfeld1992}.

The present work leaves open whether the reconstruction cover should ultimately be
interpreted as a topological cover over an emergent spacetime manifold, a site of
continuation states in the sense of Grothendieck topology, a category of
reconstruction observables, or a more general admissibility structure. The current
treatment remains intentionally agnostic on this question. A reader trained in sheaf
theory will recognize that specifying the underlying category is necessary for a
rigorous formulation, and this constitutes one of the open categorical obligations of
the framework.
\end{remark}

\subsection{Two Types of Obstruction}

Obstruction occurs when locally admissible sections fail to admit globally stable
extension across overlapping neighborhoods. Obstruction is not metaphorical: it is a
literal failure of global compatibility in the reconstruction geometry
\cite{bott1982,hatcher2002}. We distinguish two structurally distinct types whose
physical consequences differ and whose conflation would obscure both the entropy
analysis of Section~\ref{sec:entropy} and the dark matter reinterpretation of
Section~\ref{sec:darkmatter}.

\begin{definition}[Compatibility obstruction]
\label{def:compat-obs}
Let $s_\alpha \in \sheaf(U_\alpha)$ and $s_\beta \in \sheaf(U_\beta)$ be locally
admissible sections over overlapping neighborhoods. A \emph{compatibility
obstruction} exists if
\[
  \rho_{\alpha\beta}(s_\alpha) \;\neq\; \rho_{\beta\alpha}(s_\beta)
\]
on $U_\alpha \cap U_\beta$. The sections specify incompatible continuation structures
on their shared domain and cannot be globally glued.
\end{definition}

\begin{definition}[Dispersion obstruction]
\label{def:disp-obs}
A pair of sections suffers \emph{dispersion obstruction} if they agree formally on
overlapping domains --- $\rho_{\alpha\beta}(s_\alpha) = \rho_{\beta\alpha}(s_\beta)$
--- but cannot maintain stable global extension because $S$ grows faster than the
coherence propagation capacity of $\mathbf{v}$ across the cover. Formally compatible
sections therefore fail reconstruction persistence: $\mathcal{P}[s] \to 0$ under
scale extension.
\end{definition}

\begin{remark}
Not all globally inconsistent dynamics arise from local disagreement. Some arise from
uncontrolled admissibility growth under extension. Compatibility obstruction manifests
as irregular, patchy anomalies without clear scaling structure. Dispersion obstruction
manifests as smooth large-scale departure from locally stable behavior --- the
observational signature more characteristic of galactic rotation anomalies. The two
obstruction types produce distinct signatures in the entropy structure of the field
(Section~\ref{sec:entropy}) and distinct phenomenological predictions
(Appendix~\ref{app:obs}).
\end{remark}

\subsection{The Reconstruction Sheaf}

The objects defined above constitute the data of a \emph{reconstruction sheaf}
$\sheaf$ over the field domain \cite{bredon1997,kashiwara2005}: $\sheaf$ assigns to
each $U_\alpha$ the space $\sheaf(U_\alpha)$ of admissible sections of $X$ over
$U_\alpha$, together with restriction maps
\[
  \rho_{\alpha\beta} \colon \sheaf(U_\alpha) \longrightarrow \sheaf(U_\alpha \cap U_\beta)
\]
satisfying identity preservation, compatibility under composition, and the locality
and gluing axioms. Global sections of $\sheaf$ correspond to globally admissible
reconstruction modes.

\begin{remark}[Admissibility nonlocality]
\label{rem:nonlocal}
Admissibility nonlocality --- when it occurs --- is nonlocality in the indexing and
coherence structure of $\sheaf$, not in the stalk-local field dynamics. The stalks
remain local field values propagating according to local field equations. What becomes
nonlocal is the coherence condition on their assembly into globally stable sections.
This is not a superluminal signaling proposal: nonlocality appears in the geometry of
assembly rather than in microscopic propagation. See Appendix~\ref{app:locality} for
a precise separation of these notions.
\end{remark}

\subsection{Bridge to Section~\ref{sec:transition}}

Once admissibility is understood as a global extension condition rather than a local
interaction condition, it becomes natural for reconstruction dynamics to transition
between qualitatively distinct regimes as $\chi$ varies. In strongly factorizable
regimes ($\chi \ll 1$), local sections extend to global sections without reference to
cosmological boundary data. As $\chi$ approaches unity, sections begin requiring
cosmological gluing data for stable extension. The existence of a characteristic
scale --- the admissibility radius $\RA$ --- at which this transition occurs is
proposed as a consequence of the field dynamics of $X$.

% ══════════════════════════════════════════════════════════════════════════════
\section{The Admissibility Phase Transition}
\label{sec:transition}
% ══════════════════════════════════════════════════════════════════════════════

The appearance of the MOND acceleration scale \cite{milgrom1983,mcgaugh2011}
\begin{equation}
  a_0 \sim cH_0
\end{equation}
suggests that galactic dynamics are sensitive not merely to local mass distributions
but to the relation between local reconstruction stability and cosmological
admissibility structure. The coupling of galactic rotation to cosmological structure
is difficult to interpret inside purely local Newtonian dynamics --- it strongly
suggests the relevant object is not local force accumulation alone but global
admissibility structure across scales. Within RSVP, we interpret this transition not
as a modification of force law but as a \emph{breakdown of local factorization in the
reconstruction cover}, a geometric event controlled by $\chi = S/\Phi$.

\subsection{The Factorization Transition}

In regimes where $\chi \ll 1$, local admissibility neighborhoods remain strongly
factorizable. Reconstruction trajectories admit stable extension independently of
cosmological boundary data, effective dynamics reduce approximately to Newtonian
behavior, and the three entropy channels converge under coarse-graining
(Section~\ref{sec:entropy}).

As
\begin{equation}
  \chi \;=\; \frac{S}{\Phi} \;\longrightarrow\; 1,
\end{equation}
dispersion becomes comparable to continuation accessibility. Local sections cease to
remain scale-independent under extension. Admissibility neighborhoods lose strong
factorization. Reconstruction increasingly requires cosmological gluing data to
maintain persistence across the cover. Dispersion obstruction begins to dominate
compatibility structure: even formally compatible sections fail reconstruction
persistence.

\begin{definition}[Admissibility phase transition]
\label{def:phase-trans}
The \emph{admissibility phase transition} is the qualitative change in reconstruction
regime occurring when $\chi \to 1$: the transition from locally dominated
factorizable reconstruction to cosmologically coupled reconstruction requiring global
boundary data for stable extension.
\end{definition}

This is not a transition in the force law acting on matter, nor a transition in
spacetime geometry in the classical sense. It is a transition in the factorization
structure of admissible reconstruction --- a change in what kind of boundary data
local sections require for globally stable continuation. In transactional entropic
gravity \cite{kastner2024}, the force law is modified by thermodynamic corrections
applied to an otherwise standard local framework. In RSVP, the effective dynamics
change because the geometric object governing which continuations are admissible
changes.

\subsection{The Admissibility Radius}

\begin{definition}[Admissibility radius]
\label{def:RA}
The \emph{admissibility radius} $\RA$ is the largest characteristic scale over which
admissible local sections remain approximately factorizable under renormalization of
the reconstruction cover. Equivalently, $\RA$ marks the onset of
dispersion-dominated reconstruction dynamics for which stable continuation requires
cosmological extension data: a local section at scale $r$ requires cosmological
gluing data if and only if its restriction to neighborhoods approaching $\RA$ fails
to admit stable extension through the sheaf's restriction maps alone --- that is, if
$\chi$ evaluated at that scale approaches unity under the coarse-graining flow.
\end{definition}

$\RA$ is a derived quantity of the coupled dynamics of $\Phi$ and $S$, not an
external parameter. Different field configurations will generically produce different
admissibility radii.

\subsection{Dimensional Structure and Scaling Constraints}
\label{sec:dimensional}

A structural test of any geometric framework is whether the scales it identifies are
forced by its primitive objects rather than fitted to known phenomenology. We examine
the scaling dimensions of the RSVP primitives to show that the emergence of $\RA$ and
$a_0$ is not arbitrary.

Assign natural dimensions, where $[L]$ denotes length and $[T]$ time:
\begin{align*}
  [\Phi] &= L^2 T^{-2} \quad \text{(accessibility, dimension of specific energy)}, \\
  [\mathbf{v}] &= L\,T^{-1} \quad \text{(coherence-flow velocity)}, \\
  [S] &= L^2 T^{-2} \quad \text{(dispersion, matching $\Phi$ so $\chi$ is dimensionless)}, \\
  [\chi] &= 1.
\end{align*}
The control parameter $\chi = S/\Phi$ is dimensionless by construction --- the first
consistency check. The phase transition criterion $\chi \to 1$ is a pure number and
therefore scale-invariant in form, confirming that the transition is a genuine
geometric threshold rather than a unit-dependent artifact.

The admissibility radius $\RA$ has dimension $[L]$. The only combination of
fundamental constants with this dimension accessible to the reconstruction sheaf,
without introducing Newton's constant $G$ (which would bring in a mass scale absent
from the pure admissibility geometry), is $c/H_0$. The corresponding acceleration
scale satisfies
\[
  \aA \;\sim\; \frac{c^2}{\RA} \;=\; \frac{c^2}{c/H_0} \;=\; cH_0,
\]
which is dimensionally $[L\,T^{-2}]$ as required.

\begin{observation}
\label{obs:dimensional}
Within the RSVP primitive set $\{c,\, H_0,\, \chi\}$, the MOND acceleration scale
$a_0 \sim cH_0$ is the \emph{unique} dimensionally consistent admissibility
acceleration in the limit where $G$ and particle masses enter only as derived
quantities. Any alternative scaling of $\RA$ would require introducing an independent
length or time scale not present in the admissibility geometry.
Conjecture~\ref{conj:RA} is therefore dimensionally forced: the question is not
whether $\RA \propto c/H_0$ is the right form, but whether the field dynamics
uniquely select it.
\end{observation}

This argument also constrains the admissibility beta-function $\beta(\chi)$ of
Appendix~\ref{app:flow}: since $\chi$ and $\ln\ell$ are both dimensionless, $\beta$
must be a dimensionless function of $\chi$ alone at leading order. Any dimensional
parameter entering $\beta$ would break the scale-invariance of the transition
structure and produce a phase transition at a scale other than the unique admissibility
radius.

\subsection{Newtonian and Cosmologically Coupled Regimes}

\medskip
\noindent\textbf{Newtonian regime} ($r \ll \RA$, $\chi \ll 1$).\quad
Local reconstruction is strongly factorizable. Admissible sections extend stably
without cosmological boundary data. Dispersion obstruction remains bounded under
scale extension. Effective gravitational dynamics reduce to Newtonian behavior as a
consequence of local reconstruction stability rather than as a fundamental force law
\cite{verlinde2011}. Newtonian gravity therefore appears as an effective
low-dispersion reconstruction theory, analogous to the emergence of hydrodynamics
from microscopic statistical systems \cite{goldenfeld1992}.

\medskip
\noindent\textbf{Cosmologically coupled regime} ($r \sim \RA$, $\chi \sim 1$).\quad
Local factorization breaks down. Admissible sections require cosmological gluing
data. Dispersion obstruction dominates. Reconstruction persistence fails for
configurations that remain locally consistent. Effective dynamics depart from
Newtonian behavior not because the force law changes, but because the admissible
continuation set changes.

\subsection{The Central Conjecture}

\begin{conjecture}[MOND scaling from admissibility radius]
\label{conj:RA}
The weak-gradient RSVP field equations determine an asymptotic admissibility radius
satisfying
\begin{equation}
  \RA \;\propto\; \frac{c}{H_0}
\end{equation}
in low-density regimes where $\chi$ approaches unity smoothly. The corresponding
admissibility acceleration scale is then
\begin{equation}
  \aA \;\sim\; \frac{c^2}{\RA} \;\sim\; cH_0,
\end{equation}
yielding MOND-like phenomenology as a structural consequence of reconstruction
geometry.
\end{conjecture}

\begin{heuristic}
The structural plausibility of Conjecture~\ref{conj:RA} rests on the following
observation. The Hubble radius $c/H_0$ is the characteristic scale of cosmological
causal accessibility. If $\RA$ is the admissibility analog of causal accessibility,
its asymptotic scaling with $c/H_0$ would follow from the requirement that
admissibility structure and causal structure coincide asymptotically in low-density
regimes. This is a structural argument for the conjecture's reasonableness, not a
derivation.
\end{heuristic}

The MOND scale $a_0 \sim cH_0$ has been observed by several authors to couple
galactic dynamics to cosmological structure in a way that resists purely local
explanation \cite{milgrom1983,famaey2012}. Within RSVP, this coupling emerges
naturally from the admissibility framework: the transition scale is determined by the
scale at which local reconstruction loses self-sufficiency. MOND is therefore demoted
from foundational principle to phenomenological signature of a more general
reconstruction transition.

\subsection{Open Proof Obligation}

\begin{obligation}
\label{obl:RA}
Show that the coupled dynamics of $\Phi$ and $S$ in the weak-gradient limit uniquely
determine $\RA$ with asymptotic scaling $\RA \propto c/H_0$. This requires:
\emph{(i)} deriving the coarse-graining flow of $\chi$ under renormalization of the
reconstruction cover in the weak-gradient limit; \emph{(ii)} identifying the
fixed-point or threshold structure of this flow; and \emph{(iii)} showing that the
characteristic scale of the threshold is asymptotically proportional to $c/H_0$.
\end{obligation}

Conjecture~\ref{conj:RA} is falsifiable in three distinct senses: (1) if the RSVP
field equations admit a unique $\RA$ whose asymptotic scaling substantially departs
from $c/H_0$, then MOND-like phenomenology must arise from a different mechanism;
(2) if MOND-like behavior were observed in regimes where $\chi$ remains well below
unity, this would constitute evidence against the sheaf structure assumed here; and
(3) if $\RA$ can be derived but depends sensitively on unconstrained parameters in
the RSVP field equations, the conjecture is technically falsifiable but practically
underdetermined --- a third failure mode named explicitly.

\subsection{Bridge to Section~\ref{sec:entropy}}

The admissibility phase transition has consequences not only for effective dynamics
but for the entropy structure of the field. In strongly factorizable regimes, the
distinct entropy channels converge under coarse-graining because local reconstruction
is approximately scale-independent. As $\chi \to 1$, this convergence fails.
Section~\ref{sec:entropy} formalizes this consequence, deriving entropy-channel
behavior as a corollary of reconstruction geometry rather than as an independent
thermodynamic hypothesis.

% ══════════════════════════════════════════════════════════════════════════════
\section{Entropy Species as Regime Projections}
\label{sec:entropy}
% ══════════════════════════════════════════════════════════════════════════════

Most emergent gravity programs implicitly identify informational, thermodynamic,
configurational, and gravitational entropy as interchangeable quantities, treating
their equivalence as a foundational step enabling the derivation of gravitational
dynamics \cite{verlinde2011,kastner2024}. The information-thermodynamic
identification remains debated \cite{jaynes1957}, and its use as an axiom in
precisely the regimes where it is least reliable constitutes a significant structural
weakness. Within RSVP, these entropy channels are instead treated as distinct
projections of admissibility dispersion under different reconstruction covers and
coarse-grainings \cite{amari2016}. Their convergence is not assumed axiomatically
but emerges only within specific reconstruction regimes. This converts what is
typically an axiom into a prediction.

\subsection{Entropy Channels as Projections}

Given a reconstruction cover $\Ucover = \cover$, different entropy species correspond
to distinct coarse-grained projections of admissibility dispersion across
$X = (\Phi, \mathbf{v}, S)$.  We formalize these as projection functors
\[
  \Pi_{\mathrm{therm}},\quad
  \Pi_{\mathrm{info}},\quad
  \Pi_{\mathrm{grav}},\quad
  \Pi_{\mathrm{conf}}
\]
mapping admissibility structures into entropy observables associated with distinct
reconstruction regimes and coarse-graining procedures.

Informational entropy, represented by $\Pi_{\mathrm{info}}$, measures uncertainty in
continuation accessibility under local reconstruction.  It quantifies the spread of
admissible continuations available to a local section of $\Phi$ and therefore tracks
the local geometry of admissible continuation space.

Configurational entropy, represented by $\Pi_{\mathrm{conf}}$, measures admissible
arrangement multiplicity across neighboring sections of the reconstruction cover.
It characterizes the number of globally compatible assemblies obtainable from locally
admissible reconstruction data.

Thermodynamic entropy, represented by $\Pi_{\mathrm{therm}}$, measures the
coarse-grained dispersion of reconstruction trajectories under persistent exchange
processes generated by the coherence-flow field $\mathbf{v}$.  It corresponds to the
effective trajectory-volume growth observed under macroscopic reconstruction flow.

Gravitational entropy, represented by $\Pi_{\mathrm{grav}}$, measures dispersion
associated with large-scale reconstruction curvature and admissibility obstruction.
Its value depends not only on local continuation accessibility but on the obstruction
structure governing global extension of admissible sections across the reconstruction
sheaf.

None of these entropy channels is treated as fundamentally privileged within the
RSVP framework.  They correspond instead to distinct projection structures defined
over the same underlying admissibility geometry.  Their relationships are therefore
regime-dependent rather than identity-relations.  In strongly factorizable regimes,
the projection functors become approximately equivalent under admissible
coarse-graining.  Near the admissibility transition, however, the projections
separate as dispersion obstruction grows and local reconstruction loses
scale-independence.

\subsection{Entropy Convergence Under Strong Factorization}

\begin{tschema}[Entropy convergence]
\label{ts:convergence}
In reconstruction regimes satisfying $\chi \ll 1$ and admitting bounded dispersion
obstruction under renormalization of the reconstruction cover, entropy channels
converge to order~$\varepsilon$ under admissible coarse-graining transformations:
\[
  \Pitherm \;\simeq\; \Piinfo \;\simeq\; \Pigrav \;\simeq\; \Piconf
  \qquad (\chi \ll 1).
\]
Equivalently, informational, configurational, thermodynamic, and gravitational
entropy measures become approximately equivalent up to corrections of order $\chi$
and the dispersion obstruction growth rate.
\end{tschema}

When $\chi \ll 1$, local reconstruction is strongly factorizable and approximately
scale-independent. Different coarse-graining choices produce approximately the same
result because admissibility dispersion is tightly concentrated near stable
reconstruction modes --- the cover choice becomes approximately irrelevant, and the
distinct entropy projections converge. The entropy-identification moves made by
emergent gravity programs are approximately valid in this regime, but only in this
regime. The present framework derives the conditions under which they hold and
identifies them as regime-specific approximations.

\subsection{Entropy Divergence Under Dispersion Obstruction}

As $\chi \to 1$, dispersion obstruction begins to dominate. Different coarse-graining
choices produce systematically different results. Informational entropy, tracking
local continuation accessibility, may remain bounded while thermodynamic entropy
grows. Gravitational entropy, tracking obstruction contributions, may diverge from
configurational entropy as global extension failures accumulate. Functorially:
\[
  \Pitherm \;\not\simeq\; \Pigrav
  \qquad (\chi \to 1).
\]

\begin{remark}
Persistent divergence between entropy species is not an annoying inconsistency to be
resolved by a better identification. It is a physical signal: evidence of obstruction
growth within the reconstruction geometry itself. When entropy channels diverge
persistently in a physical system, the RSVP framework predicts this as a signature
of the system operating near or beyond the admissibility phase transition.
\end{remark}

\subsection{Commutator Structure of Entropy Projections}
\label{sec:commutator}

The entropy projection functors $\Pi_i$ are not interchangeable under coarse-graining.
This non-commutativity is a structural feature of the framework, not a defect.

Define the coarse-graining operator $\Rcoarse$ acting on reconstruction covers as in
Appendix~\ref{app:flow}. The commutator of a projection with coarse-graining measures
whether applying the projection before or after scale-averaging produces the same
result:
\[
  [\Pi_i,\, \Rcoarse] \;:=\; \Pi_i \circ \Rcoarse \;-\; \Rcoarse \circ \Pi_i.
\]
In strongly factorizable regimes ($\chi \ll 1$), the projections approximately commute
with coarse-graining:
\[
  [\Pi_i,\, \Rcoarse] \;\approx\; 0 \qquad (\chi \ll 1).
\]
This is precisely the regime in which Theorem-Schema~\ref{ts:convergence} holds:
entropy species converge because different orderings of projection and coarse-graining
yield the same result.

Near the admissibility transition ($\chi \to 1$), the commutator becomes nontrivial:
\[
  [\Pi_{\mathrm{therm}},\, \Rcoarse] \;\not\approx\; 0 \qquad (\chi \to 1).
\]
The thermodynamic projection and gravitational projection respond differently to
coarse-graining near the transition because they track different aspects of
admissibility dispersion --- local trajectory averaging versus large-scale obstruction
growth, respectively. Their divergence is therefore not a failure of the framework
but a prediction: it identifies the admissibility phase transition as the point at
which entropy-channel ordering becomes physically meaningful.

\begin{remark}
The non-commutativity $[\Pi_i, \Rcoarse] \neq 0$ near $\chi \to 1$ is the entropy-
theoretic manifestation of the same breakdown captured geometrically by dispersion
obstruction (Definition~\ref{def:disp-obs}). Both say the same thing from different
angles: coarse-graining and local projection cease to commute when reconstruction
factorization fails. This convergence of geometric and entropic descriptions onto the
same transition criterion is a structural consistency check.
\end{remark}

\subsection{The Kastner--Schlatter Contrast}

Kastner and Schlatter \cite{kastner2024} identify Shannon information entropy with
thermodynamic entropy as a foundational step, using this identification to drive the
derivation of gravitational dynamics. Within RSVP, this identification is
approximately valid when $\chi \ll 1$ (Theorem-Schema~\ref{ts:convergence}) and
fails as a foundational assumption precisely in the regimes of greatest physical
interest: the weak-gradient, low-density regimes where MOND-like behavior appears
and $\chi$ approaches unity. Founding the derivation on entropy identification in
those regimes imports an approximation as an axiom exactly where the approximation
breaks down.

\subsection{Bridge to Section~\ref{sec:darkmatter}}

The existence of persistent entropy-channel divergence without local dynamical
inconsistency suggests that globally anomalous galactic behavior need not arise from
missing mass distributions. Recall Remark~\ref{rem:cover-derived}: the reconstruction
cover is a derived object of the field geometry, not a background against which the
field is defined. Anomalies in the global extension structure of the cover are
anomalies in the geometry itself, not in the matter content of a pre-given space.

% ══════════════════════════════════════════════════════════════════════════════
\section{Dark Matter as Reconstruction Obstruction}
\label{sec:darkmatter}
% ══════════════════════════════════════════════════════════════════════════════

\subsection{The Observational Problem}

The observational signatures attributed to dark matter are well established
empirically: flat galactic rotation curves \cite{mcgaugh2011,famaey2012}, weak
gravitational lensing excesses, cluster dynamics requiring additional gravitational
sources, and large-scale structure formation rates exceeding baryonic predictions.

Standard cosmology interprets these anomalies as evidence for additional unseen mass
distributions embedded within a fixed spacetime background. Within RSVP, the
anomalies are instead interpreted as \emph{signatures of obstruction in the global
extension structure of the reconstruction cover}. As established in
Remark~\ref{rem:cover-derived}: the reconstruction cover is a derived object of the
field geometry, not a background against which the field is defined. Anomalies in
the global extension structure of the cover are anomalies in the geometry itself.

\subsection{Compatibility and Dispersion Obstruction in Galactic Dynamics}

Compatibility obstruction (Definition~\ref{def:compat-obs}) would manifest as local
dynamical inconsistency --- irregular, patchy anomalies without clear scaling
structure. The observational signatures of galactic anomalies do not have this
character. Galaxies are locally stable and internally coherent; their stellar and gas
dynamics are well described by standard hydrodynamics and Newtonian gravity at small
radii.

The anomaly appears in the \emph{persistence} of those dynamics under large-scale
reconstruction. This is structurally closer to dispersion obstruction
(Definition~\ref{def:disp-obs}). Local sections remain formally compatible across
neighboring regions; the failure is not disagreement but persistence failure under
scale extension: as $\chi$ approaches unity near $\RA$, local reconstruction dynamics
lose scale-independence and can no longer be extended stably without cosmological
boundary data. The point is not that galaxies behave inconsistently. The point is
that their local reconstruction dynamics fail to remain self-sufficient under
cosmological extension.

\subsection{Rotation Curves as Persistence Failure}

In strongly factorizable regimes ($\chi \ll 1$), orbital dynamics are governed
predominantly by local continuation accessibility, and the admissible trajectory set
is closed under local extension --- the regime in which Newtonian dynamics holds as
an emergent approximation. As dispersion obstruction grows and $\chi$ approaches
unity near $\RA$, local continuation sets cease to remain closed under scale
extension. The admissible trajectory structure at galactic radii approaching $\RA$
is no longer determined by local field data alone --- it requires cosmological gluing
data for stable continuation. The effective admissible trajectory set therefore
departs from Newtonian closure even in the absence of additional local mass.

The transition criterion $\chi \to 1$ marks the onset of this regime --- the same
criterion governing the admissibility phase transition
(Section~\ref{sec:transition}) and the entropy-channel divergence
(Section~\ref{sec:entropy}). The unification of all three phenomena under the same
control parameter is the framework's strongest structural feature.

\begin{observation}
\label{obs:phase-transition-central}
\emph{The anomaly is a geometric phase transition.} Flat rotation curves, MOND-like
scaling, and entropy-channel divergence are manifestations of the same underlying
transition in the factorization structure of admissible reconstruction, controlled by
$\chi \to 1$ near $\RA$.
\end{observation}

\subsection{Weak Lensing and Extension Geometry}

Within RSVP, photon propagation probes the continuation accessibility geometry of
$\Phi$ rather than merely local mass density. In strongly factorizable regimes,
continuation accessibility tracks local mass distributions closely and lensing
reduces approximately to standard gravitational lensing. In regimes approaching the
admissibility phase transition, the continuation accessibility geometry becomes
influenced by obstruction structure in the reconstruction cover.

\begin{obligation}
\label{obl:lensing}
Compute the effective lensing potential from the obstruction curvature of $\sheaf$
in the weak-gradient limit and compare with observed convergence maps, including
Bullet Cluster phenomenology.
\end{obligation}

The framework identifies this as a possible reinterpretation of lensing anomalies in
terms of obstruction curvature. Whether the resulting phenomenology quantitatively
reproduces observed lensing profiles remains an open problem. The Bullet Cluster
remains a decisive empirical challenge: the framework must reproduce the observed
lensing structure through obstruction curvature of the reconstruction geometry. The
restraint here is deliberate; this target is stated honestly rather than
rhetorically dismissed.

\subsection{Entropy Divergence as Observational Signal}

Regions dominated by dispersion obstruction should exhibit increasing divergence
between informational, thermodynamic, configurational, and gravitational entropy
measures. Gravitational entropy, tracking obstruction contributions, should depart
from thermodynamic entropy of local baryonic processes in exactly the regimes where
rotation curve anomalies appear. Standard $\Lambda$CDM does not predict
entropy-channel divergence as a geometric signal; RSVP predicts that such divergence
is a structural signature of the admissibility phase transition.

\begin{obligation}
\label{obl:entropy-signal}
Derive explicit expressions for the gravitational entropy contribution of
reconstruction obstruction in the weak-gradient RSVP limit and determine whether
predicted divergence onset correlates with observed rotation curve transition radii.
\end{obligation}

\subsection{Situating the Reinterpretation}

The RSVP reinterpretation differs structurally from all three major competing
programs. Particle dark matter models preserve local factorization while adding
unseen sources. Entropic gravity programs \cite{verlinde2017,padmanabhan2010}
preserve entropy identification while modifying effective dynamics. MOND and TeVeS
\cite{milgrom1983,bekenstein2004} modify the force law through fitted interpolation
functions. RSVP instead interprets galactic anomalies as signatures of breakdown in
scale-independent reconstruction itself --- the onset of dispersion-dominated
extension dynamics controlled by $\chi$ approaching unity near $\RA$
(Observation~\ref{obs:phase-transition-central}).

% ══════════════════════════════════════════════════════════════════════════════
\section{Comparison with Neighboring Programs}
\label{sec:comparison}
% ══════════════════════════════════════════════════════════════════════════════

\subsection{Organizing Principle}

The most useful axis for comparing competing approaches is derivation architecture:
what object is treated as primitive, what undergoes transition, what role entropy
plays, and what counts as an anomaly.

\begin{table}[ht]
\centering
\renewcommand{\arraystretch}{1.35}
\begin{tabular}{@{}p{3.2cm}p{2.8cm}p{2.8cm}p{2.5cm}p{2.5cm}@{}}
\toprule
\textbf{Framework} & \textbf{Primitive} & \textbf{Anomaly} & \textbf{Entropy} & \textbf{Locality} \\
\midrule
$\Lambda$CDM & Matter source & Missing source & Identified & Local field \\
MOND/TeVeS & Modified force & Wrong dynamics & Secondary & Modified local \\
Verlinde & Entropy equiv.\ & Incomplete thermo.\ & Axiom & Holographic \\
Trans.\ entropic & Exchange completion & Incomplete thermo.\ & Interchangeable & Local exchange \\
RSVP & Admissibility geometry & Phase transition & Regime projection & Sheaf-indexed \\
\bottomrule
\end{tabular}
\caption{Comparative taxonomy organized by derivation architecture.}
\label{tab:comparison}
\end{table}

\subsection{Standard Cosmology ($\Lambda$CDM)}

$\Lambda$CDM treats matter source distributions within a fixed pseudo-Riemannian
background as primitive. What RSVP interprets as obstruction structure in a derived
cover, $\Lambda$CDM interprets as source deficits within a fixed background. These
are not empirically equivalent restatements: they predict different entropy-channel
signatures, different lensing phenomenology in obstruction-dominated regimes, and
different structure formation dynamics. RSVP constrains the anomaly structure through
$\chi$ and $\RA$; the phenomenology is not freely adjustable but is derived from
field dynamics.

\subsection{MOND and TeVeS}

MOND \cite{milgrom1983} reproduces flat rotation curves and the baryonic
Tully-Fisher relation \cite{mcgaugh2011} with remarkable precision using a single
additional parameter. Within the present framework, MOND is demoted from foundational
principle to phenomenological signature. The MOND acceleration scale $a_0 \sim cH_0$
is not a new fundamental constant but a derived geometric scale --- the acceleration
associated with the onset of cosmologically coupled reconstruction at $\RA$.
MOND-like dynamics appear when $\chi \to 1$ near $\RA$, not because the force law
has been modified but because the admissible continuation set has changed.

The observed robustness of the baryonic Tully-Fisher relation \cite{mcgaugh2011}
across galaxy types, which is difficult to understand in $\Lambda$CDM and built into
MOND by construction, would in RSVP reflect universality of the admissibility phase
transition under coarse-graining --- an emergent property of the reconstruction
geometry rather than a tuned interpolation (see Appendix~\ref{app:rg}).

\subsection{Verlinde's Emergent Gravity}

Verlinde \cite{verlinde2011,verlinde2017} begins with entropy identification and
derives effective geometry. RSVP begins with admissibility geometry and derives the
conditions under which entropy species approximately converge
(Theorem-Schema~\ref{ts:convergence}). Verlinde's framework inherits the
entropy-identification problem \cite{jaynes1957} as a foundational commitment. RSVP
converts this axiom into a theorem-schema valid in strongly factorizable regimes and
failing in precisely the regimes --- $\chi \to 1$ --- where anomalous dynamics appear.

\subsection{Transactional Entropic Gravity}

Kastner and Schlatter \cite{kastner2024} provide the most immediately adjacent
comparison. Both programs reject fundamental spacetime; both treat geometric structure
as emergent. The divergence is in derivation order and primitive object.

Transactional entropic gravity treats exchange process completion as primitive ---
spacetime links are generated by completed transactions between irreducible emitters
and absorbers. RSVP treats admissibility geometry as primitive --- emitters,
absorbers, and exchange processes are stabilized roles within the reconstruction
sheaf.

Three structural consequences follow: (1) the transactional MOND derivation depends
on entropy-species identification at specific steps; the RSVP conjecture depends on a
single geometric condition ($\chi \to 1$ at $\RA$) controlling all phenomena
simultaneously; (2) transactional completion is strictly weaker than sheaf-compatible
global extension --- a transaction can complete locally while generating globally
inadmissible dynamics; (3) entropy-species interchangeability appears in RSVP as an
approximate consequence of strongly factorizable reconstruction, valid in the
Newtonian regime and failing in the MOND regime.

RSVP does not currently outperform transactional entropic gravity quantitatively. The
advantage claimed is architectural: fewer primitive objects controlling more phenomena
through a single transition structure, with entropy convergence as a derived
consequence.

\subsection{The Inversion as Theoretical Commitment}

The organizing inversion --- admissibility geometry generates exchange structure, not
vice versa --- reflects a specific theoretical bet about where the explanatory
primitives of gravitational physics are located. If the bet is correct, existing
phenomenology should emerge as special cases of reconstruction geometry. If wrong,
the transactional program has the correct derivation order and RSVP's sheaf structure
is descriptive rather than explanatory. The paper makes the question precise and
identifies the geometric objects on which it turns.

% ══════════════════════════════════════════════════════════════════════════════
\section{Minimal Assumptions and Irreducible Commitments}
\label{sec:assumptions}
% ══════════════════════════════════════════════════════════════════════════════

The RSVP framework has grown sufficiently rich that it is important to distinguish
which assumptions are load-bearing from which are illustrative scaffolding. This
section identifies the irreducible commitments of the theory --- the propositions
that, if abandoned, would dissolve the framework rather than merely simplify it.

\medskip
\noindent\textbf{(I) Admissibility geometry is primary.}
The reconstruction cover $\Ucover$, admissibility classes, and the sheaf $\sheaf$
are not derived from a more primitive physical ontology of particles, fields on fixed
spacetime, or exchange processes. They are the foundational objects from which
physical roles (emitter, absorber, particle, clock, force) emerge as stabilized
reconstruction modes. This commitment is what distinguishes RSVP from transactional
and holographic programs.

\medskip
\noindent\textbf{(II) The reconstruction cover is derived, not background.}
There is no pre-given manifold against which the field evolves. The topology of
admissibility neighborhoods is determined dynamically by $\chi = S/\Phi$. Classical
spacetime, to the extent it appears, is a stable low-$\chi$ regime of the cover, not
its presupposition (Remark~\ref{rem:cover-derived}).

\medskip
\noindent\textbf{(III) Coarse-graining modifies factorization structure.}
Under renormalization of the reconstruction cover, what changes is not merely
effective couplings (as in standard EFT) but the factorization structure of admissible
reconstruction itself. This is what allows admissibility nonlocality to appear without
modifying local propagation (Appendix~\ref{app:locality}).

\medskip
\noindent\textbf{(IV) Entropy convergence is a regime consequence, not an axiom.}
The approximate identification of entropy species holds when $\chi \ll 1$ and breaks
down as $\chi \to 1$. This is a prediction, not a foundational assumption
(Theorem-Schema~\ref{ts:convergence}, Section~\ref{sec:commutator}).

\medskip
\noindent\textbf{(V) Obstruction is dynamical, not merely topological.}
The distinction between compatibility and dispersion obstruction
(Definitions~\ref{def:compat-obs},~\ref{def:disp-obs}) is essential: obstruction
classes can grow under coarse-graining even when they are trivializable at a fixed
scale. This dynamical obstruction is what the framework uses to interpret galactic
anomalies.

\medskip
Everything else in the paper --- the candidate dynamical equations of
Appendix~\ref{app:dynamics}, the specific form of the beta-function in
Appendix~\ref{app:flow}, the detailed structure of the Lean schemas in
Appendix~\ref{app:implementation} --- is illustrative rather than irreducible. Those
elements demonstrate that the five commitments above are compatible with coherent
dynamics and formal implementation, but they do not exhaust the space of theories
satisfying (I)--(V).

\begin{remark}
The minimality of commitments (I)--(V) is itself a structural feature. A framework
whose central claims reduce to five propositions, from which all major results are
derived, is more constrained than one requiring independent axioms for each
phenomenon. The recurrence of $\chi$, $\RA$, and the obstruction classes across
Sections~\ref{sec:geometry}--\ref{sec:darkmatter} is evidence that the commitments
are not merely compatible but actively generative.
\end{remark}

% ══════════════════════════════════════════════════════════════════════════════
\section{The Great Inversion and the Ontological Status of Physical Law}
\label{sec:inversion}
% ══════════════════════════════════════════════════════════════════════════════

A persistent ambiguity in contemporary approaches to gravity concerns the ontological
ordering between microscopic processes and macroscopic geometry. In most emergent
gravity programs, local interactions, thermodynamic exchanges, or entanglement
structures are treated as fundamental, with spacetime geometry reconstructed as a
large-scale byproduct. The RSVP framework reverses this ordering. Admissibility
geometry is not derived from physical exchange; rather, physical exchange appears only
insofar as it is permitted as a stable reconstruction mode within a pre-existing
geometric admissibility structure.

This inversion eliminates a hidden circularity present in many bottom-up
constructions. If local exchanges are taken as primitive, one must still specify which
exchanges are allowed to persist across time, and under what conditions they remain
stable. These persistence conditions are often introduced implicitly through
thermodynamic or informational assumptions. In the RSVP framework, such assumptions
are elevated to the level of geometric constraint. The admissibility structure
determines which continuations are permitted, and physical processes are those
continuations that achieve stable realization.

Under this interpretation, particles are not fundamental objects in the usual sense.
They are persistent resonant modes within the admissibility geometry: localized
configurations where continuation accessibility and coherence flow align sufficiently
to maintain stability under extension. The apparent solidity of matter reflects a
regime of strong reconstruction persistence rather than a primitive ontology.
Observational anomalies are therefore not interpreted as evidence for missing entities
within a fixed background, but as indications that the reconstruction structure itself
has entered a regime where its factorization properties fail.

Traditional physical theories typically begin by specifying primitive entities,
fields, background geometry, and dynamical laws governing their interaction. The RSVP
framework instead begins from admissible continuation, persistence under extension,
compatibility of reconstruction, and obstruction constraints governing global
coherence. This constitutes a shift from ontology-first physics toward
constraint-first physics. Physical law becomes increasingly interpretable as
constraint structure governing globally persistent continuation rather than as a
description of material substance.

The recurring mathematical objects of the framework --- $\chi$, $\RA$, $\Omega_\ell$,
$\sheaf$, and the entropy projection functors --- all function primarily as
constraints on admissible extension rather than as descriptions of substance. Sheaf
theory is fundamentally concerned with conditions under which local consistency
extends to global coherence, and that is precisely the structural role admissibility
plays within RSVP\@. The framework therefore increasingly resembles a generalized
consistency geometry: its central objects are not material constituents but
admissibility relations, extension constraints, persistence criteria, and obstruction
structures governing which reconstructions can stabilize across scale.

% ══════════════════════════════════════════════════════════════════════════════
\section{The Control Parameter and the Emergence of Physical Regimes}
\label{sec:chi-interpretation}
% ══════════════════════════════════════════════════════════════════════════════

The dynamical behavior of the framework is governed not by the individual fields
$\Phi$, $\mathbf{v}$, or $S$ in isolation, but by the dimensionless ratio $\chi =
S/\Phi$. This ratio functions as a regime-defining parameter that determines the
structural character of admissible reconstruction. Its role is analogous to that of
an order parameter, but it governs the factorization properties of reconstruction
rather than the phase of a material system.

When $\chi \ll 1$, admissibility dispersion is negligible relative to continuation
accessibility. Admissible trajectories concentrate around stable modes, and local
sections extend without requiring information from larger scales. The resulting
structure is strongly factorizable: local dynamics close under extension, and the
behavior of physical systems can be determined from local data alone. Classical
Newtonian dynamics and general relativity emerge as effective descriptions precisely
in this regime. Their apparent universality reflects the dominance of accessibility
over dispersion, not a fundamental status.

As $\chi$ approaches unity, this structure degrades. The dispersion of admissible
continuations becomes comparable to the accessibility itself, and local reconstruction
loses closure under extension. The breakdown is structural rather than dynamical: the
governing equations do not change, but the conditions under which local solutions
extend consistently are no longer satisfied. This transition explains the empirical
breakdown of classical dynamics at galactic scales without requiring any modification
of the underlying force laws. The laws remain intact; the regime in which they
provide a closed description has ceased to apply.

This reframes the MOND phenomenology. Rather than representing a modification of
Newtonian dynamics, the MOND transition marks the boundary between regimes of
admissibility geometry. The apparent change in behavior reflects a transition in the
structure of admissible continuations rather than a change in the dynamical equations
themselves. The acceleration scale $a_0 \sim cH_0$ is not a new constant of nature
but the scale associated with the breakdown of local reconstruction closure --- a
geometric threshold determined by the admissibility structure of $\Phi$ and $S$.

% ══════════════════════════════════════════════════════════════════════════════
\section{Observable Proxies for Reconstruction Geometry}
\label{sec:observables}
% ══════════════════════════════════════════════════════════════════════════════

A persistent difficulty for any reconstruction-theoretic framework is the distinction
between formally meaningful geometric objects and operationally measurable quantities.
The RSVP program introduces objects such as admissibility dispersion, persistence
obstruction, and reconstruction curvature whose mathematical role is structurally
clear while their observational realization remains incomplete. This section
distinguishes sharply between primitive reconstruction objects and their possible
empirical proxies, and in doing so identifies the operational bridge that a successful
completion of the theory must eventually construct.

The coupled field $X = (\Phi, \mathbf{v}, S)$ is not assumed to be directly
observable. Observations access only coarse-grained effective projections of
admissibility structure. Consequently, the relevant empirical question is not whether
$\Phi$, $\chi$, or $\Omega_\ell$ are directly measurable, but whether stable
observational signatures exist whose behavior systematically tracks their
reconstruction dynamics.

The control parameter $\chi = S/\Phi$ is interpreted operationally as a measure of
reconstruction instability under scale extension. Since neither $S$ nor $\Phi$ is
independently observable in the present formulation, $\chi$ must be inferred
indirectly through signatures of factorization breakdown. Candidate observational
proxies include the onset radii of rotational persistence anomalies, mismatch between
baryonic predictions and large-scale orbital stability, environment-sensitive
deviations from MOND interpolation structure, and entropy-channel divergence between
thermodynamic and gravitational reconstruction measures.

The obstruction functional $\Omega_\ell[s]$ should not be interpreted as directly
measurable cohomological data. Instead, $\Omega_\ell$ represents an effective
obstruction intensity whose observable manifestations may include weak lensing
residuals after baryonic subtraction, persistence instability under environmental
perturbation, anomalous scaling behavior near galactic transition radii, and
cluster-scale reconstruction mismatch.

The framework therefore predicts not merely modified dynamics but systematic
correlations between $\chi$, $\Omega_\ell$, entropy-channel divergence, and
persistence stability. The existence or nonexistence of these correlations constitutes
an empirical test of the reconstruction interpretation.

\begin{remark}
The present framework does not yet provide a completed observational dictionary
mapping reconstruction geometry to astrophysical observables. The purpose of this
section is narrower: to identify the operational bridge that a successful completion
of the theory must eventually construct. A framework that can specify what would count
as empirical support or falsification, even before the measurement is possible, is
structurally more mature than one whose observational content remains entirely
implicit.
\end{remark}

% ══════════════════════════════════════════════════════════════════════════════
\section{Admissibility Selection Principles and Forbidden Continuation Classes}
\label{sec:forbidden}
% ══════════════════════════════════════════════════════════════════════════════

One of the principal risks of generalized reconstruction frameworks is excessive
permissiveness. If every sufficiently smooth local continuation can be rendered
globally admissible through appropriate reinterpretation of the reconstruction cover,
then obstruction loses predictive force and the framework degenerates into descriptive
topology. The RSVP program therefore requires a distinction between formal
constructibility and persistent admissibility.

Not every mathematically definable reconstruction corresponds to a physically
persistent reconstruction mode. A continuation structure must satisfy increasingly
restrictive admissibility criteria under extension and coarse-graining flow. We
therefore distinguish the following hierarchy: constructible configurations form the
broadest class; locally admissible configurations must additionally satisfy bounded
perturbation stability; globally extendable configurations must also cohere across the
full reconstruction cover without compatibility obstruction; and persistent
configurations must furthermore survive renormalization of the cover without
dispersion obstruction growth exceeding coherence transport capacity.

\begin{definition}[Forbidden continuation class]
A continuation class is \emph{forbidden} if no admissible coarse-graining flow
preserves reconstruction persistence under extension beyond finite scale.
Equivalently, every admissible extension flow eventually drives the persistence
functional $\mathcal{P}[s]$ to zero.
\end{definition}

Forbidden continuation classes represent geometries of reconstruction that cannot
stabilize as globally persistent physical regimes. The framework predicts that certain
classes of admissibility structure are dynamically excluded regardless of their local
formal consistency. Candidate forbidden classes include runaway dispersion geometries
with unbounded entropy-channel divergence, obstruction cascades whose growth exceeds
coherence transport capacity, non-factorizable continuation structures lacking stable
attractor behavior, and reconstruction flows possessing no bounded admissibility fixed
point.

This exclusion structure is essential. Mature physical theories derive much of their
predictive strength not from explanatory flexibility but from impossibility structure.
The RSVP program therefore increasingly shifts from ontology-generation toward
admissibility exclusion. The transition from descriptive admissibility geometry to
exclusion geometry marks the most important maturation step available to the
framework: a successful reconstruction theory must eventually identify not only what
can persist, but what cannot.

% ══════════════════════════════════════════════════════════════════════════════
\section{Known Structural Risks and Failure Modes}
\label{sec:risks}
% ══════════════════════════════════════════════════════════════════════════════

The RSVP framework remains incomplete and structurally vulnerable in several
important respects. Naming these risks explicitly is important both for conceptual
clarity and for preventing the framework from expanding into unconstrained
abstraction. A framework that understands its own fragility structure is structurally
more mature than one that presents only its successes.

The most significant risk is \emph{semantic overextension}. The framework employs a
common vocabulary --- admissibility, reconstruction, persistence, obstruction,
continuation, and projection --- across gravitational, thermodynamic, informational,
and geometric domains. This unification is potentially powerful, but creates a risk
that the same concepts become sufficiently elastic to redescribe phenomena without
generating unique constraints. The framework therefore lives or dies by whether the
same small set of structures --- $\chi$, $\RA$, $\Omega_\ell$, and the reconstruction
sheaf --- repeatedly generate nontrivial predictions across independent domains rather
than functioning as a universal explanatory solvent.

A related risk is \emph{universality inflation}. The paper repeatedly suggests that
distinct microscopic reconstruction dynamics may flow toward common large-scale
admissibility behavior under coarse-graining. While this possibility is motivated by
ordinary renormalization theory, the present work does not yet derive the relevant
admissibility operator algebra, identify the relevant versus irrelevant operators
governing admissibility flow, or prove the existence of universality classes for
reconstruction dynamics. The universality hypothesis therefore remains speculative,
and a critic asking why admissibility flow should possess universality classes at all
has not yet been given a satisfying answer.

The framework also faces \emph{observational opacity}. Many central objects remain
operationally indirect, and the theory currently lacks a completed observational
dictionary relating admissibility quantities to uniquely measurable astrophysical
observables. Without such a dictionary, the framework risks remaining geometrically
expressive but empirically underdetermined. The proxies identified in
Section~\ref{sec:observables} are a first step toward resolving this, but they remain
candidate mappings rather than completed operational definitions.

\emph{Coarse-graining ambiguity} represents a fourth structural vulnerability. The
admissibility coarse-graining operator $\Rcoarse$ is not yet rigorously defined. The
framework draws analogies from Wilsonian renormalization while modifying the
factorization structure of admissible continuation itself, but has not yet specified
what degrees of freedom are integrated out, how reconstruction neighborhoods evolve
under flow, or what invariants are preserved. The admissibility coarse-graining
operation is not assumed to coincide with ordinary Wilsonian integration over momentum
shells --- it acts instead on the factorization structure of admissible continuation
itself --- but until the operator is rigorously constructed, the beta-function
formalism of Appendix~\ref{app:flow} remains heuristic.

\emph{Insufficient exclusion structure} is a fifth risk. If admissibility conditions
remain too permissive, the framework risks collapsing into generalized reconstruction
semantics in which every phenomenon can be redescribed in admissibility language after
the fact. The introduction of forbidden continuation classes in
Section~\ref{sec:forbidden} is intended as a step toward resolving this, but the
relevant exclusion theorems remain undeveloped.

Finally, the \emph{categorical status} of the reconstruction sheaf remains
unresolved. It is currently unclear whether the admissibility cover should ultimately
be interpreted as a topological cover over emergent spacetime, a continuation-state
site, a reconstruction category, or a more general higher-categorical admissibility
structure. The present framework remains intentionally agnostic on this question, but
a reader trained in sheaf theory will eventually demand an answer to the question of
what category these sheaves actually live over.

\begin{remark}
These risks should not be interpreted merely as weaknesses. They identify the exact
structural obligations whose resolution would determine whether RSVP matures into a
genuine physical reconstruction theory or remains a descriptive geometric framework.
A framework that can name the conditions under which it would fail is structurally
more credible than one whose explanatory architecture can indefinitely absorb anomaly
through reinterpretation.
\end{remark}

% ══════════════════════════════════════════════════════════════════════════════
\section{Criteria for a Successful RSVP Completion}
\label{sec:completion}
% ══════════════════════════════════════════════════════════════════════════════

Speculative frameworks often fail to define what successful completion would actually
look like, allowing the program to defer resolution indefinitely into abstraction. The
RSVP program therefore requires explicit completion criteria that give the framework a
visible finish line.

A successful reconstruction-theoretic completion of RSVP would minimally require a
derived admissibility beta-function from explicit RSVP field dynamics; a rigorous
definition of admissibility coarse-graining specifying what degrees of freedom are
integrated out and what invariants are preserved; operational observables
corresponding to $\chi$ and obstruction structure; a derivation of $\RA \propto
c/H_0$ from weak-gradient reconstruction flow without parameter fitting; recovery of
Einstein-like behavior in strongly factorizable regimes as an emergent consequence
rather than a postulate; quantitative lensing predictions from obstruction curvature
capable of distinguishing the framework from dark matter models on Bullet Cluster
scales; proof or refutation of entropy-channel convergence under strong factorization;
classification of forbidden continuation classes with accompanying exclusion theorems;
and demonstration that the framework generates unique empirical constraints rather
than post-hoc reinterpretations.

More broadly, the framework must eventually answer three foundational questions. Which
reconstruction geometries persist? Which reconstruction geometries are forbidden?
Which observable structures uniquely distinguish admissibility geometry from
ordinary matter-source explanations? The success or failure of the RSVP program
should ultimately be judged not by the breadth of its interpretive vocabulary but by
its ability to constrain continuation structure, generate exclusion geometry, and
produce empirically distinguishable predictions.

\begin{remark}
A framework that can state the conditions under which it would fail is structurally
more mature than one whose explanatory architecture can indefinitely absorb anomaly
through reinterpretation. The completion criteria above are not aspirational rhetoric
but genuine falsifiability conditions: if the framework cannot eventually satisfy
them, the reconstruction interpretation should be abandoned in favor of programs with
stronger exclusion structure.
\end{remark}

% ══════════════════════════════════════════════════════════════════════════════
\section{Conclusion}
\label{sec:conclusion}
% ══════════════════════════════════════════════════════════════════════════════

\subsection{The Central Reinterpretation}

The present work develops a reinterpretation of gravitational anomalies organized
around a single foundational inversion. Existing emergent gravity programs treat
physical exchange, thermodynamic bookkeeping, or information flow as primitive and
derive spacetime geometry from them. RSVP treats admissibility geometry as primary
and interprets persistent physical exchange as a stabilized reconstruction mode.
Emitters, absorbers, particles, force laws, and thermodynamic entropy are not inputs
to the framework but outputs --- roles stabilized by recursive cross-scale coherence
of the coupled field $\Xfield$ within its reconstruction sheaf.

The framework's explanatory structure is organized by a small set of recurring
mathematical objects: the coupled field $X$; the control parameter $\chi = S/\Phi$;
the reconstruction cover $\Ucover$; the admissibility radius $\RA$; the
compatibility and dispersion obstruction classes; and the entropy projection functors
$\Pitherm, \Piinfo, \Pigrav, \Piconf$. The control parameter $\chi$ governs
simultaneously: Newtonian gravitational stability, the MOND-like transition,
entropy-channel convergence, and galactic obstruction phenomenology. This unification
under a single order parameter is the strongest evidence of internal coherence.

\subsection{What the Paper Claims and Does Not Claim}

The paper does not present: a completed derivation of MOND from RSVP field equations;
a quantitative alternative to $\Lambda$CDM; or a solved lensing model capable of
reproducing Bullet Cluster phenomenology.

The paper does present: a reconstruction-geometric reinterpretation of gravitational
anomalies organized by a stable conceptual hierarchy; formal definitions of
admissibility geometry, reconstruction cover, and obstruction types; a constrained
and explicitly falsifiable conjecture (Conjecture~\ref{conj:RA}) with three named
failure modes; Theorem-Schema~\ref{ts:convergence} converting an axiom of neighboring
programs into a derived consequence; and a comparative taxonomy organized by
derivation architecture (Table~\ref{tab:comparison}).

\begin{quote}
The present work does not derive MOND or dark matter phenomenology from completed
RSVP field equations; it identifies the geometric objects from which such derivations
would have to proceed.
\end{quote}

\subsection{Open Problems as Evidence of Structural Coherence}

Most speculative frameworks accumulate explanatory flexibility. The present framework
instead generates constrained open problems whose resolution would confirm or falsify
its central commitments:

\begin{enumerate}[label=(\arabic*)]
  \item \textbf{Admissibility radius derivation} (Obligation~\ref{obl:RA}). Show
    that the weak-gradient RSVP equations uniquely determine $\RA \propto c/H_0$.
  \item \textbf{Universality class structure.} Determine whether distinct admissibility
    geometries flow toward the same weak-gradient reconstruction behavior under
    coarse-graining, making MOND phenomenology geometrically necessary.
  \item \textbf{Lensing derivation target} (Obligation~\ref{obl:lensing}). Compute
    the effective lensing potential from obstruction curvature of $\sheaf$ and
    compare with observed convergence maps.
  \item \textbf{Entropy divergence signature} (Obligation~\ref{obl:entropy-signal}).
    Identify observational contexts where gravitational entropy departs measurably
    from thermodynamic entropy of local baryonic processes.
  \item \textbf{Obstruction flow analysis.} Determine whether compatibility and
    dispersion obstruction produce distinguishable galactic signatures varying
    systematically with galaxy morphology, surface brightness, or environment.
\end{enumerate}

These problems are obligations generated by the framework's own internal structure.
Their existence as constrained derivation targets --- rather than free parameters or
post-hoc interpretive choices --- is evidence that the framework has crossed the
threshold where speculative architecture begins behaving as a research program. A
framework that knows the shape of its own incompleteness is structurally more mature
than one that knows only the shape of its successes.

\subsection{Final Philosophical Position}

If the RSVP program is correct, then gravity, inertia, dark matter phenomenology,
and entropy equivalence are not independent physical ingredients assembled into a
cosmological model. They are stable regimes and failure modes of reconstruction
geometry --- consequences of the factorization structure of admissible continuation
under a derived cover whose topology is controlled by a single dimensionless
parameter.

Whether this inversion reflects physical reality remains unresolved. The present work
aims only to make the inversion mathematically explicit, empirically constrained, and
structurally falsifiable.

% ══════════════════════════════════════════════════════════════════════════════
\appendix
% ══════════════════════════════════════════════════════════════════════════════

\section{Formal Structure of the Reconstruction Sheaf}
\label{app:sheaf}

This appendix develops the formal sheaf-theoretic structure underlying the
reconstruction framework. The purpose is not to introduce new physical claims but to
make explicit the mathematical objects already used operationally throughout
Sections~\ref{sec:geometry}--\ref{sec:darkmatter}.

\subsection{Presheaf and Sheaf Structure}

Let $\recon$ denote the reconstruction domain and $\Ucover = \cover$ a cover by
admissibility neighborhoods. The \emph{reconstruction presheaf} $\sheaf$ assigns to
each $U_\alpha$ a set $\sheaf(U_\alpha)$ of admissible sections of $X$ over
$U_\alpha$, together with restriction maps
\[
  \rho_{V}^{U} \colon \sheaf(U) \longrightarrow \sheaf(V), \qquad V \subseteq U,
\]
satisfying $\rho_{U}^{U} = \mathrm{id}$ and $\rho_{W}^{V} \circ \rho_{V}^{U} =
\rho_{W}^{U}$ for $W \subseteq V \subseteq U$.

$\sheaf$ is a \emph{sheaf} if it satisfies the following for every open cover
$\{U_i\}$ of any $U$:
\begin{enumerate}[label=(\roman*)]
  \item \textit{Locality:} if $\rho_{U_i}^{U}(s) = \rho_{U_i}^{U}(t)$ for all $i$,
    then $s = t$.
  \item \textit{Gluing:} if $\rho_{U_i \cap U_j}^{U_i}(s_i) =
    \rho_{U_i \cap U_j}^{U_j}(s_j)$ for all $i, j$, then there exists $s \in
    \sheaf(U)$ with $\rho_{U_i}^{U}(s) = s_i$.
\end{enumerate}

Compatibility obstruction (Definition~\ref{def:compat-obs}) corresponds to failure
of condition (ii): sections cannot be glued because they disagree on overlaps.
Dispersion obstruction (Definition~\ref{def:disp-obs}) is subtler: sections satisfy
(ii) formally but the glued section fails reconstruction persistence under
coarse-graining because admissibility dispersion increases faster than coherence
propagation.

\subsection{Admissible Sections and Persistence}

An admissible local section is not merely a field restriction but a structured pair
consisting of: a local restriction of the coupled field
\[
  X = (\Phi,\mathbf{v},S)
\]
to an admissibility neighborhood $U_\alpha$, together with a continuation structure
specifying the admissible extension behavior of reconstruction trajectories
originating within $U_\alpha$.

Formally, a section
\[
  s_\alpha \in \sheaf(U_\alpha)
\]
is admissible only if two conditions are simultaneously satisfied.  First, local
perturbations of the section must preserve bounded continuation behavior, so that
small deformations of initial reconstruction data do not generate uncontrolled
trajectory divergence.  Second, admissible extensions of the section must remain
stable under renormalization of the reconstruction cover.  The section must therefore
persist not only locally but under scale-dependent restructuring of admissibility
neighborhoods induced by coarse-graining flow.

This second condition distinguishes reconstruction persistence from ordinary local
field consistency.  A section may satisfy all local field equations and remain
formally compatible on overlaps while nevertheless failing persistence under
extension because admissibility dispersion grows faster than coherence propagation
under cover renormalization.

A global reconstruction mode corresponds to a global section
\[
  s \in \sheaf(\recon)
\]
whose restrictions agree coherently across the full reconstruction cover and whose
continuation structure remains stable under admissible coarse-graining flow.
Persistence is therefore a multiscale extension property of the reconstruction
sheaf rather than merely a statement of local dynamical consistency.

\subsection{Persistence Functional}

We define a \emph{persistence functional}
\[
  \mathcal{P}[s] \;\in\; [0,1]
\]
measuring stability of admissible extensions under renormalization of the
reconstruction cover. A section satisfies reconstruction persistence if
$\mathcal{P}[s]$ remains bounded away from zero under scale extension. Dispersion
obstruction occurs when $\mathcal{P}[s] \to 0$ under coarse-graining, even for
sections satisfying the standard gluing conditions.

\subsection{Dynamical Cover and Gluing}

Unlike ordinary differential-geometric frameworks in which the manifold is fixed and
fields evolve upon it, the RSVP framework treats the effective reconstruction
topology as dependent upon admissibility flow. In RSVP, gluing conditions are
dynamically weighted by admissibility dispersion and continuation accessibility
rather than determined solely by overlap consistency. The resulting structure combines:
local field dynamics, scale-dependent admissibility flow, and global extension
stability into a single reconstruction structure.

The following commutative diagram summarizes the restriction structure for a triple
cover:

\[
\begin{tikzcd}
\sheaf(U_\alpha) \arrow[r, "\rho_{\alpha\beta}"] \arrow[dr, "\rho_{\alpha\gamma}"']
  & \sheaf(U_\alpha \cap U_\beta) \arrow[d, "\rho_{\alpha\beta,\gamma}"] \\
  & \sheaf(U_\alpha \cap U_\beta \cap U_\gamma)
\end{tikzcd}
\]

Consistency of this diagram for all triples is the cocycle condition whose failure
generates compatibility obstructions classified by $\Hcech^1(\Ucover, \sheaf)$.

\subsection{Constraint Propagation and Reconstruction Stability}
\label{app:constraint-prop}

Admissibility is defined locally but must persist globally. A key missing link between
the sheaf language and the field dynamics is an account of how local admissibility
violations propagate under coarse-graining. We sketch the required structure here.

Define a \emph{reconstruction constraint functional}
\[
  \mathcal{C}[X] \;=\; 0
\]
whose vanishing characterizes admissible field configurations. In the candidate
dynamics of Appendix~\ref{app:dynamics}, this constraint is approximately realized
when $\chi \ll 1$ and reconstruction trajectories are strongly localized near stable
modes. The constraint is violated as $\chi \to 1$.

Define \emph{propagation operators} at scale $\ell$:
\[
  \mathcal{T}_\ell \colon \sheaf(U_\ell) \;\longrightarrow\; \sheaf(U_{\lambda\ell}),
  \qquad \lambda > 1,
\]
describing how admissible sections at scale $\ell$ map to sections at the coarser
scale $\lambda\ell$. Reconstruction stability requires that the constraint propagates
covariantly:
\[
  \mathcal{T}_\ell\bigl(\mathcal{C}[X]\bigr) \;=\; 0
\]
whenever $\mathcal{C}[X] = 0$ at the finer scale. Failure of this condition --- a
section satisfying $\mathcal{C} = 0$ locally but $\mathcal{T}_\ell(\mathcal{C}) \neq
0$ at the coarser scale --- is precisely dispersion obstruction: local admissibility
that does not propagate stably under renormalization of the cover.

The propagation operators $\{\mathcal{T}_\ell\}$ form a directed system. Their
composition law
\[
  \mathcal{T}_{\lambda\ell} \circ \mathcal{T}_\ell \;=\; \mathcal{T}_{\lambda^2\ell}
\]
is a consistency condition analogous to the cocycle condition for restriction maps. A
complete theory of constraint propagation would derive the $\mathcal{T}_\ell$ from the
field equations of Appendix~\ref{app:dynamics} and show that constraint violation
first appears at $\ell \sim \RA$ --- formalizing the admissibility phase transition as
a constraint-propagation threshold rather than merely a qualitative change in
reconstruction regime.

\begin{obligation}
\label{obl:constraint}
Derive the propagation operators $\mathcal{T}_\ell$ from the coupled dynamics of
$\Phi$, $\mathbf{v}$, and $S$ and show that $\mathcal{T}_\ell(\mathcal{C}) \neq 0$
first appears at $\ell \sim \RA \propto c/H_0$ in weak-gradient regimes. This would
convert the admissibility phase transition from a geometric threshold into a
propagation-failure theorem.
\end{obligation}
\section{Coarse-Graining Flow and the Admissibility Order Parameter}
\label{app:flow}

The main text introduced $\chi = S/\Phi$ as the central control parameter. This
appendix develops a schematic coarse-graining framework and clarifies the sense in
which the admissibility phase transition may possess universality structure analogous
to ordinary critical phenomena \cite{wilson1974,goldenfeld1992}.

\subsection{Scale-Dependent Fields and Order Parameter}

Let $\Ucover_\ell$ denote the reconstruction cover at coarse-graining scale $\ell$.
Integrating out short-scale reconstruction modes modifies both continuation
accessibility and admissibility dispersion:
\[
  \Phi \;\longrightarrow\; \Phi_\ell, \qquad S \;\longrightarrow\; S_\ell.
\]
The scale-dependent order parameter is
\[
  \chi_\ell \;=\; \frac{S_\ell}{\Phi_\ell}.
\]
The Newtonian regime ($\chi_\ell \ll 1$) corresponds to strongly factorizable
reconstruction; the admissibility transition occurs when $\chi_\ell \to 1$.

\subsection{Admissibility Beta-Function}

The coarse-graining operation in RSVP is not assumed to coincide with ordinary
Wilsonian integration over momentum shells. Wilsonian coarse-graining integrates out
short-wavelength field modes while preserving the background spacetime manifold.
Informational coarse-graining aggregates over microstates while preserving statistical
ensemble structure. Admissibility-cover renormalization instead acts on the
factorization structure of admissible continuation itself: it merges reconstruction
neighborhoods and asks whether the resulting coarser cover still admits globally
persistent sections. These three operations are structurally distinct and should not
be conflated. The admissibility beta-function introduced below governs the third
operation, not the first two.

The coarse-graining flow of the reconstruction geometry may be represented
schematically by the scale evolution equation
\begin{equation}
  \label{eq:beta}
  \frac{d\chi_\ell}{d\ln \ell}
  \;=\;
  \beta(\chi_\ell),
\end{equation}
where
\[
  \chi_\ell
  =
  \left.\frac{S}{\Phi}\right|_{\ell}
\]
denotes the scale-dependent control parameter evaluated at coarse-graining scale
$\ell$, and $\beta$ is an \emph{admissibility beta-function} governing the
renormalization flow of reconstruction stability.

The function $\beta$ is not specified phenomenologically but constrained by the
structural requirements of admissible reconstruction developed throughout the main
text.  In strongly factorizable regimes,
\[
  \chi \ll 1,
\]
the framework requires
\[
  \beta(\chi) < 0,
\]
so that admissibility dispersion decreases under coarse-graining flow and local
reconstruction remains stable under scale extension.  This corresponds to the
Newtonian regime in which reconstruction neighborhoods retain approximate
factorization and admissible sections persist without cosmological gluing data.

As $\chi$ approaches unity, the flow enters a crossover region in which dispersion
obstruction begins competing with coherence propagation.  The qualitative structure
of admissible extension changes: local sections cease to remain scale-independent,
and persistence under renormalization becomes increasingly sensitive to global
reconstruction geometry.  The admissibility phase transition introduced in the main
text corresponds to this change in flow structure near
\[
  \chi \sim 1.
\]

Beyond the transition, the framework admits two structurally distinct possibilities.
Either the flow approaches a new large-scale stable reconstruction regime governed
by cosmologically coupled admissibility structure, or local factorization ceases to
remain dynamically stable altogether and reconstruction persistence fails beyond
finite scale.  Which possibility is realized depends on the weak-gradient dynamics
of the RSVP field equations and constitutes part of the open derivation program
identified in Obligation~\ref{obl:RA}.

The admissibility beta-function therefore plays a role analogous to renormalization
group flow equations in statistical field theory.  It governs not merely numerical
coupling evolution but qualitative transitions in the factorization structure of
admissible reconstruction itself.

\subsection{Universality Classes}

Suppose multiple microscopic RSVP field dynamics produce distinct local
reconstruction behavior while nevertheless generating the same asymptotic
coarse-graining flow for $\chi_\ell$. Then the weak-gradient galactic regime would
exhibit universal phenomenology independent of microscopic reconstruction details ---
analogous to universality of critical exponents in condensed matter systems
\cite{goldenfeld1992}. This possibility is particularly important for understanding
the empirical robustness of MOND-like scaling relations \cite{mcgaugh2011,famaey2012}
across diverse galactic systems. Within RSVP, such robustness would emerge as a
universality property of admissibility flow near the reconstruction transition rather
than a tuned interpolation.

\subsection{Admissibility Radius from Flow}

The admissibility radius $\RA$ is the characteristic scale at which $\chi_\ell \to
1$ under the coarse-graining flow. Obligation~\ref{obl:RA} is therefore equivalent
to showing that flow equation~\eqref{eq:beta}, derived from the weak-gradient RSVP
field equations, admits a threshold at scale $\ell^*$ satisfying $e^{\ell^*} \propto
c/H_0$.

\subsection{Attractor Structure and Stability of Newtonian Reconstruction}
\label{app:attractor}

The Newtonian regime ($\chi \ll 1$) is not merely a parameter region; it should be
an \emph{attractor} of the admissibility flow. This subsection formalizes that
requirement.

If the beta-function satisfies $\beta(\chi) < 0$ for $\chi \ll 1$, then the fixed
point $\chi^* = 0$ is stable: perturbations of $\chi$ away from zero are driven back
under coarse-graining. Newtonian reconstruction is therefore an infrared attractor of
the admissibility flow in strongly factorizable regimes. This is the geometric
meaning of Newtonian gravity's empirical robustness at small scales.

The admissibility phase transition at $\chi \sim 1$ corresponds to the boundary of
this attractor basin. For $\chi < 1$, the flow returns to the Newtonian attractor.
For $\chi \geq 1$, the flow either enters a new cosmologically coupled basin or
diverges, depending on the sign of $\beta$ beyond the transition.

A minimal consistent beta-function with this attractor structure takes the form:
\[
  \beta(\chi) \;=\; -\mu\,\chi\,(1 - \chi) \;+\; \mathcal{O}(\chi^3)
\]
for some $\mu > 0$. This function satisfies:
\begin{itemize}
  \item $\beta(0) = 0$: the Newtonian fixed point $\chi^* = 0$ is a fixed point.
  \item $\beta'(0) = -\mu < 0$: the fixed point is stable (attractor).
  \item $\beta(1) = 0$: the transition point $\chi = 1$ is an unstable fixed point
    (the admissibility phase boundary).
  \item $\beta(\chi) < 0$ for $0 < \chi < 1$: the flow drives $\chi$ toward zero
    in the Newtonian regime.
\end{itemize}
This schematic form is consistent with all qualitative requirements established in
the main text. Deriving the actual beta-function from the RSVP field equations ---
including the value of $\mu$ and the structure of higher-order terms --- constitutes
part of Obligation~\ref{obl:RA}.

\begin{remark}
If the attractor structure is correct, the universality of Newtonian gravity at small
scales is not a fundamental fact but a consequence of reconstruction flow. Deviations
from Newtonian behavior at large scales (galactic rotation anomalies) are then not
anomalies requiring new matter but signatures of crossing the attractor boundary ---
the transition from the $\chi^* = 0$ basin to the cosmologically coupled regime.
\end{remark}
\section{Candidate Dynamical Equations for the RSVP Fields}
\label{app:dynamics}

This appendix sketches candidate dynamical equations for the coupled fields
$\Xfield$. These equations are heuristic and incomplete --- candidate structures
consistent with the admissibility framework, not finalized physical laws. The purpose
is architectural: to show that the reconstruction-geometric interpretation can
plausibly be associated with a coherent field-dynamical system.

\subsection{Continuation Accessibility Field $\Phi$}

A minimal admissibility-flow equation for $\Phi$:
\begin{equation}
  \label{eq:Phi}
  \partial_t\Phi \;=\; D_\Phi\nabla^2\Phi - \lambda_\Phi \Phi S
    + \gamma_\Phi \nabla\cdot\mathbf{v} + N_\Phi[\Phi,\mathbf{v},S]
\end{equation}
where $D_\Phi$ is an accessibility diffusion coefficient, $\lambda_\Phi$ governs
suppression of continuation accessibility by dispersion growth ($-\lambda_\Phi \Phi
S$ captures the central geometric tension: increasing $S$ suppresses stable $\Phi$),
$\gamma_\Phi$ couples accessibility to coherence flow divergence, and $N_\Phi$
denotes higher-order nonlinear terms.

\subsection{Coherence Flow Field $\mathbf{v}$}

\begin{equation}
  \label{eq:v}
  \partial_t\mathbf{v} \;=\; D_v\nabla^2\mathbf{v} - (\mathbf{v}\cdot\nabla)\mathbf{v}
    + \alpha_v\nabla\Phi - \lambda_v S\mathbf{v} + N_v[\Phi,\mathbf{v},S]
\end{equation}
The gradient term $\alpha_v\nabla\Phi$ drives coherence flow toward high continuation
accessibility. The suppression term $-\lambda_v S\mathbf{v}$ models degradation of
coherent transport under increasing dispersion. The nonlinear transport term resembles
fluid advection but its interpretation is geometric: reconstruction coherence
propagates directionally through admissibility structure.

\subsection{Entropy-Dispersion Field $S$}

\begin{equation}
  \label{eq:S}
  \partial_t S \;=\; D_S\nabla^2 S + \sigma_S |\nabla\Phi|^2
    - \mu_S S\Phi - \nabla\cdot(S\mathbf{v}) + N_S[\Phi,\mathbf{v},S]
\end{equation}
The source term $\sigma_S |\nabla\Phi|^2$ generates dispersion through accessibility
curvature gradients. The suppression term $-\mu_S S\Phi$ stabilizes dispersion in
regions of high continuation accessibility. The advection term $-\nabla\cdot(S\mathbf{v})$
transports dispersion through coherence flow.

\subsection{Control Parameter Dynamics}

Together, equations \eqref{eq:Phi}--\eqref{eq:S} determine the evolution of the
control parameter:
\[
  \partial_t\chi \;=\; \frac{\Phi\,\partial_t S - S\,\partial_t\Phi}{\Phi^2}.
\]
Equation~\eqref{eq:S} is of particular interest: the source term $\sigma_S(1 -
S/\Phi)S = \sigma_S(1 - \chi)\chi\Phi$ drives $\chi$ toward unity when $\chi < 1$
and away when $\chi > 1$, producing the admissibility transition at $\chi = 1$ as a
dynamical fixed point.

\subsection{Epistemic Status}

These equations are schematic. Many alternative systems may generate the same
admissibility phenomenology; if the universality hypothesis of Appendix~\ref{app:flow}
is correct, weak-gradient reconstruction behavior may be largely insensitive to
microscopic dynamical details. A completed RSVP field theory would require: a
variational principle, covariant generalization, explicit reconstruction observables,
well-posed initial value structure, and quantitative comparison with data.

% ── App D ─────────────────────────────────────────────────────────────────────
\section{Obstruction Cohomology and Global Reconstruction Failure}
\label{app:cohomology}

This appendix sketches a cohomological interpretation of reconstruction obstruction
\cite{bott1982,hatcher2002,kashiwara2005}. The purpose is to demonstrate
mathematical directionality, not claim a finished cohomological theory.

\subsection{\v{C}ech Cohomology of the Reconstruction Cover}

Let $\Ccech^\bullet(\Ucover, \sheaf)$ denote the \v{C}ech cochain complex associated
with $\Ucover$ and $\sheaf$. A family of local sections $\{s_\alpha\}$ defines a
$0$-cochain. A compatibility obstruction corresponds to a non-trivial $1$-cocycle
$g_{\alpha\beta} \in \Ccech^1(\Ucover, \sheaf)$ measuring the failure of sections to
agree on overlaps. Globally admissible reconstruction requires trivialization of
this cocycle:

\begin{definition}[Compatibility obstruction class]
\label{def:obs-class}
A \emph{compatibility obstruction class} is a nontrivial element of
\[
  \Hcech^1(\Ucover,\sheaf),
\]
corresponding to the impossibility of globally gluing locally admissible sections.
Vanishing $\Hcech^1$ indicates that all compatible local sections admit global
extension without compatibility obstruction.
\end{definition}

\subsection{Two Obstruction Notions}

Within RSVP, obstruction possesses an additional dynamical feature absent from
ordinary static sheaf theory: the reconstruction cover evolves under admissibility
coarse-graining. This motivates two notions:

\begin{definition}[Static obstruction]
A \emph{static obstruction} is an ordinary compatibility obstruction at a fixed
reconstruction scale --- a nontrivial class in $\Hcech^1(\Ucover_\ell, \sheaf)$ for
some fixed $\ell$.
\end{definition}

\begin{definition}[Persistence obstruction]
A \emph{persistence obstruction} occurs when an obstruction class fails to be
trivialized under coarse-graining of the reconstruction cover, even if it is
trivializable at some initial scale. Formally, there exists a coarse-graining flow
$\Rcoarse$ such that the induced map
\[
  \Rcoarse^* \colon \Hcech^1(\Ucover_{\lambda_1}, \sheaf) \longrightarrow
  \Hcech^1(\Ucover_{\lambda_2}, \sheaf), \quad \lambda_2 > \lambda_1,
\]
does not preserve triviality.
\end{definition}

\subsection{Scale-Dependent Obstruction Functional}

Introduce a scale-dependent obstruction functional:
\[
  \Omega_\ell[s] \;\in\; \mathbb{R}_{\geq 0}
\]
whose flow under coarse-graining satisfies:
\begin{equation}
  \frac{d\Omega_\ell}{d\ln\ell} \;=\; \Gamma\!\left(\chi_\ell, \Omega_\ell\right)
\end{equation}
where $\chi_\ell = S_\ell/\Phi_\ell$. In strongly factorizable regimes ($\chi_\ell
\ll 1$), obstruction flow remains bounded and global reconstruction modes persist
stably. Near the admissibility transition ($\chi_\ell \to 1$), dispersion obstruction
may drive $\Omega_\ell \to \infty$ even in formally compatible local systems.

\subsection{Galactic Anomalies as Cohomological Failure}

The galactic anomaly reinterpretation of Section~\ref{sec:darkmatter} may be
reformulated cohomologically: flat rotation curves and lensing anomalies correspond
to growth of persistence obstruction in the global reconstruction geometry rather
than to missing matter sources. This clarifies why anomalies appear primarily in
large-scale extension behavior rather than in local dynamical inconsistency.

A completed reconstruction-cohomology program would require: explicit
admissibility-chain complexes, construction of scale-dependent obstruction operators,
definition of persistence cohomology under admissibility renormalization, and
derivation of observational signatures associated with obstruction growth.

% ── App E ─────────────────────────────────────────────────────────────────────
\section{Relation to Renormalization and Effective Field Theory}
\label{app:rg}

This appendix clarifies the relation between RSVP's reconstruction geometry and
ordinary Wilsonian effective field theory \cite{wilson1974,goldenfeld1992}. The
purpose is not to identify RSVP with standard EFT but to explain where the analogy
is structurally useful and where the two frameworks diverge.

\subsection{Where the Analogy Holds}

In ordinary EFT, short-scale degrees of freedom are integrated out, producing
scale-dependent effective couplings governing long-scale behavior. Both frameworks
study persistence of effective structure under scale transformation. The analogy
motivates the beta-function framework of Appendix~\ref{app:flow} and the universality
class hypothesis.

\subsection{Where the Analogy Breaks}

In ordinary EFT, the background geometry remains fixed throughout renormalization;
only couplings evolve. In RSVP, coarse-graining modifies not merely effective
couplings but the factorization structure of admissible reconstruction itself. This
is the formal content of Remark~\ref{rem:cover-derived}: the reconstruction cover is
a derived object of the field geometry, not a background.

Ordinary EFT assumes that locality survives coarse-graining even as effective
couplings change. RSVP instead allows the factorization structure supporting effective
locality to become unstable. The notion of universality acquires a modified meaning:
not that many microscopic systems flow toward the same macroscopic effective
behavior, but that distinct admissibility geometries may flow toward the same
reconstruction-transition structure under coarse-graining.

\subsection{Newtonian Gravity as Effective Reconstruction Theory}

Newtonian dynamics emerges not as a fundamental interaction but as an effective
low-dispersion reconstruction theory --- analogous to hydrodynamics emerging from
microscopic statistical systems \cite{goldenfeld1992}. Hydrodynamic laws are not
fundamental microscopic truths; they are stable large-scale descriptions arising from
strong universality under coarse-graining. Within RSVP, Newtonian dynamics plays the
same role.

\subsection{Entropy Equivalence as Emergent Universality}

Entropy convergence (Theorem-Schema~\ref{ts:convergence}) is interpreted as an
effective coarse-grained phenomenon analogous to emergent thermodynamic equations of
state. Entropy equivalence is not foundational; it is a low-dispersion universality
property of strongly factorizable reconstruction. Divergence of entropy channels near
the admissibility transition is then unsurprising: universality breaks down because
reconstruction factorization itself weakens under renormalization flow.

% ── App F ─────────────────────────────────────────────────────────────────────
\section{Observational and Experimental Targets}
\label{app:obs}

This appendix organizes the empirical targets generated by the framework. It does not
claim observational confirmation; its purpose is operational --- to identify concrete
empirical obligations and failure conditions.

\subsection{Rotation Curve Transition Structure}

The framework predicts that transition from Newtonian to MOND-like behavior should
correlate with reconstruction-factorization breakdown ($\chi \to 1$) rather than
solely with local acceleration magnitude. This differs subtly from ordinary MOND:
RSVP predicts that transition behavior should correlate not only with acceleration
but with indicators of reconstruction persistence, including local coherence
structure, environmental reconstruction density, and scale-dependent entropy
divergence. Low-surface-brightness galaxies are especially important because they
probe weak-gradient regimes where $\chi \to 1$ over large spatial domains
\cite{famaey2012}.

\subsection{Entropy-Channel Divergence}

The most distinctive RSVP prediction: in strongly factorizable regimes ($\chi \ll
1$), entropy channels converge approximately. Near the admissibility transition
($\chi \sim 1$), the framework predicts increasing divergence between channels due to
dispersion obstruction growth. Standard $\Lambda$CDM and ordinary entropic gravity do
not predict entropy-channel divergence as a geometric signal. Operationally, galactic
regions exhibiting strong rotational anomalies should exhibit measurable mismatch
between local thermodynamic entropy indicators and large-scale gravitational
reconstruction structure.

\subsection{Weak Lensing Signatures}

Within RSVP, lensing anomalies should correlate with obstruction curvature of the
reconstruction cover rather than uniquely with inferred dark matter distributions.
Systems with similar baryonic matter distributions but different reconstruction
persistence structure may exhibit different lensing behavior. The Bullet Cluster
remains a decisive challenge: the framework must reproduce its lensing structure
through obstruction curvature geometry (Obligation~\ref{obl:lensing}).

\subsection{Cosmological Structure Formation}

Large-scale structure formation depends not merely on matter density perturbations
but on admissibility persistence under cosmological extension. The admissibility
radius may evolve dynamically during cosmological history as reconstruction dispersion
changes under expansion-like coarse-graining flow. These possibilities remain
speculative; the framework currently lacks quantitative cosmological simulations
capable of testing them.

\subsection{Observational Failure Modes}

The framework generates clear failure conditions: (1) if weak-gradient galactic
behavior cannot be associated with coherent admissibility-transition structure, the
geometric reinterpretation fails; (2) if cluster-scale lensing cannot be reproduced
through obstruction curvature without reintroducing hidden matter, the framework
loses explanatory advantage over $\Lambda$CDM; (3) if entropy-channel divergence
fails to correlate with anomalous galactic regimes, one of RSVP's most distinctive
predictions is falsified; (4) if $\RA$ cannot be derived without arbitrary parameter
tuning, the central conjecture becomes structurally underdetermined.

% ── App G ─────────────────────────────────────────────────────────────────────
\section{Locality, Causality, and Admissibility Nonlocality}
\label{app:locality}

The framework employs concepts such as cosmological gluing data, admissibility
nonlocality, and global extension dependence. These risk misunderstanding if
interpreted as superluminal influence or hidden signaling. This appendix specifies
precisely what type of nonlocality RSVP permits and what type it explicitly rejects.

\subsection{The Central Distinction}

The framework does \emph{not} propose superluminal communication, violation of
relativistic causal structure, or hidden instantaneous signaling between local field
degrees of freedom.  The central distinction introduced by RSVP is instead the
distinction between locality of propagation and locality of admissibility.

The stalk-local dynamics of the reconstruction sheaf remain local.  Local field
values propagate according to local dynamical equations governing the coupled field
\[
  X = (\Phi,\mathbf{v},S),
\]
and admissible perturbations evolve continuously through neighboring reconstruction
regions.  In this sense, RSVP preserves ordinary local propagation structure at the
level of stalk dynamics.

What becomes nonlocal is not propagation itself but the coherence conditions
governing extension and assembly of local sections into globally persistent
reconstruction modes.  A local section may satisfy all local dynamical equations
while nevertheless failing global admissibility because its continuation structure
cannot be extended coherently across the reconstruction cover under renormalization
flow.

Most ordinary field theories implicitly identify locality of propagation with
locality of admissibility.  RSVP separates these notions.  Locality of propagation
is a property of stalk dynamics.  Locality of admissibility is a property of global
extension structure in the reconstruction sheaf.

Admissibility nonlocality therefore appears not in microscopic signaling but in the
indexing, restriction, and gluing structure of admissible reconstruction itself.
The dependence on cosmological gluing data discussed throughout the main text is not
a proposal for instantaneous causal influence across spacetime.  It is a statement
that persistence of admissible continuation may depend on global reconstruction
coherence conditions extending beyond the local neighborhood of a given section.

This distinction becomes especially important near the admissibility transition
\[
  \chi = \frac{S}{\Phi} \sim 1,
\]
where local reconstruction ceases to remain strongly factorizable under
coarse-graining.  In strongly factorizable regimes,
\[
  \chi \ll 1,
\]
local admissibility appears effectively self-contained because global extension
conditions reduce approximately to local continuation structure.  Near the
transition, however, persistence becomes increasingly sensitive to large-scale
reconstruction geometry even while local propagation remains relativistically local.

RSVP therefore weakens locality of admissibility without abandoning locality of
propagation.  The framework modifies the geometry of globally admissible extension,
not the local causal structure of field evolution.

\subsection{Admissibility Nonlocality Defined}

Admissibility nonlocality exists in: the indexing structure of the cover, the gluing
conditions, and the persistence structure of extension under renormalization. It does
\emph{not} exist in the stalk-local dynamical equations themselves. Nonlocality
appears in the geometry of assembly rather than in microscopic propagation.

A local section may satisfy all local dynamical equations while nevertheless failing
to admit globally stable extension because admissibility persistence depends upon
large-scale coherence of the reconstruction cover. The dependence is
\emph{extension dependence}, not \emph{dynamical signaling dependence}.

\subsection{Analogy with Gauge Theory and General Relativity}

This distinction is analogous in spirit to global constraint structure in gauge
theory and general relativity. In electromagnetism, local electric-field measurements
are constrained globally by Gauss-law structure. In general relativity, local
geometry is constrained by global consistency conditions on spacetime slicing. RSVP
generalizes this: local reconstruction trajectories remain locally dynamical, but
admissibility of their persistent extension depends upon global reconstruction
geometry.

\subsection{Interpretation of ``Cosmological Gluing Data''}

The term does not refer to signals propagating instantaneously from cosmological
distances. It refers to the fact that globally persistent continuation may require
compatibility with large-scale reconstruction structure. The admissible continuation
set of a local trajectory depends upon the global extension geometry of the
reconstruction cover --- an extension condition, not a propagation condition.

Near the admissibility transition ($\chi \to 1$), local factorization weakens. Local
continuation trajectories cease to remain closed under extension independently of
cosmological admissibility structure. The resulting dependence may appear nonlocal
from the perspective of ordinary local field theory even though no superluminal
propagation occurs.

\subsection{Formal Statement}

Admissibility nonlocality occurs in the indexing and coherence structure of $\sheaf$
rather than in the stalk-local field dynamics. A completed RSVP field theory would
require rigorous proof that admissibility nonlocality remains compatible with
relativistic causal structure under all admissible reconstruction flows. The present
appendix establishes only the conceptual separation necessary for the framework to
avoid immediate conflict with ordinary causal locality. This separation follows
directly from the central inversion: persistent physical exchange is generated by
admissibility geometry rather than vice versa. If admissibility geometry is primary,
then locality of admissibility and locality of propagation need not coincide.

% ── App H ─────────────────────────────────────────────────────────────────────
\section{Formal Implementation: Grammar, Typed Structures, and Proof Schemas}
\label{app:implementation}

This appendix outlines possible formal implementation directions for the RSVP
admissibility framework. The guiding principle is that the framework should
eventually be expressible as a typed formal system in which admissibility,
obstruction, entropy-channel convergence, and reconstruction persistence can be
checked, simulated, or proven under specified assumptions.

\subsection{BNF Grammar for Reconstruction Objects}

A minimal formal language must represent fields, regions, covers, sections,
restriction maps, obstruction classes, entropy channels, and scale transformations.

\begin{lstlisting}[caption={BNF grammar for RSVP reconstruction theory}]
<Theory>     ::= <FieldDecl> <CoverDecl> <SheafDecl>
                 <DynamicsDecl> <ObstructionDecl> <EntropyDecl>

<FieldDecl>  ::= "field" <Name> "=" "(" <Phi> "," <v> "," <S> ")"
<Phi>        ::= "Phi" ":" <Domain> "->" Real
<v>          ::= "v"   ":" <Domain> "->" Vector
<S>          ::= "S"   ":" <Domain> "->" Real

<Domain>     ::= "ReconDomain" | <Region>
<Region>     ::= "U" <Index> | <Region> "cap" <Region>
                            | <Region> "cup" <Region>

<CoverDecl>  ::= "cover" <Name> "=" "{" <RegionList> "}"

<SheafDecl>  ::= "sheaf" <Name> ":" <CoverName> "->" <SectionSpace>
<SectionSpace> ::= "Sections" "(" <FieldName> "," <Region> ")"

<RestrMap>   ::= "restrict" <Section> "from" <Region> "to" <Region>

<Admiss>     ::= "admissible" "(" <Section> ")"
               | "stable"     "(" <Section> ")"
               | "persistent" "(" <Section> ")"

<Chi>        ::= "chi" "=" <S> "/" <Phi>

<ObsDecl>    ::= "obstruction" <ObsType> "(" <Section> "," <Section> ")"
<ObsType>    ::= "compatibility" | "dispersion"

<EntropyDecl>::= "entropy" <EntType> "(" <Section> ")"
<EntType>    ::= "informational" | "configurational"
               | "thermodynamic" | "gravitational"

<DynDecl>    ::= "flow" <Name> ":" <Scale> "->" <TheoryState>

<Formula>    ::= <Admiss> | <ChiRel> | <EntRel> | <ObsRel>
               | <Formula> "and" <Formula>
               | <Formula> "implies" <Formula>

<ChiRel>     ::= "chi << 1" | "chi ~ 1" | "chi -> 1"
<EntRel>     ::= "converges" "(" <EntTypeList> ")"
               | "diverges"  "(" <EntTypeList> ")"
<ObsRel>     ::= "bounded" "(" <ObsType> ")"
               | "grows"   "(" <ObsType> ")"
\end{lstlisting}

\subsection{BNF for Epistemic Status of Claims}

\begin{lstlisting}[caption={BNF grammar for claim classification}]
<Claim> ::= <Definition>
          | <Conjecture>
          | <TheoremSchema>
          | <ProofObligation>
          | <PhenomInterp>

<Definition>     ::= "define"    <Object> "as" <FormalExpr>
<Conjecture>     ::= "conjecture" <Statement> "under" <Assumptions>
<TheoremSchema>  ::= "theorem-schema" <Statement> "provided" <Conditions>
<ProofObligation>::= "prove"     <Statement> "from" <Defs> "and" <Dynamics>
<PhenomInterp>   ::= "interpret" <Observation> "as" <FrameworkObject>

<Object>    ::= "AdmissibilityRadius"    | "CompatObstruction"
              | "DispersionObstruction"  | "EntropyChannel"
              | "ReconstructionCover"    | "AdmissPhaseTransition"

<Relation>  ::= "determines" | "converges_to" | "diverges_from"
              | "scales_as"  | "obstructs"    | "extends_to"
\end{lstlisting}

This grammar prevents category errors: ``MOND is derived'' is not yet a theorem; it
is a conjectural derivation target. ``The admissibility radius is the relevant
geometric object'' is a definitional and structural claim. ``Entropy channels
converge under strong factorization'' is a theorem-schema requiring proof under
specified dynamics.

\subsection{Typed Structures in Lean Style}

\begin{lstlisting}[caption={Lean-style typed structures for RSVP objects}]
structure ReconDomain where
  points : Type

structure Region (R : ReconDomain) where
  carrier : Set R.points

structure RSVPField (R : ReconDomain) where
  Phi : R.points -> Real
  v   : R.points -> Vector
  S   : R.points -> Real

-- Control parameter (well-defined when Phi > 0)
def PositiveAccess {R : ReconDomain}
    (X : RSVPField R) : Prop :=
  (*$\forall$*) x : R.points, X.Phi x > 0

def chi {R : ReconDomain}
    (X : RSVPField R) (h : PositiveAccess X)
    (x : R.points) : Real :=
  X.S x / X.Phi x

-- Admissibility predicates
structure LocalSection {R : ReconDomain}
    (X : RSVPField R) (U : Region R) where
  field_restr : U.carrier -> RSVPField R
  continuation_data : Type
  bounded_stability : Prop

def StronglyFactorizable {R : ReconDomain}
    (X : RSVPField R) (h : PositiveAccess X)
    (U : Region R) ((*$\delta$*) : Real) : Prop :=
  (*$\forall$*) x (*$\in$*) U.carrier, chi X h x < (*$\delta$*)

def NearTransition {R : ReconDomain}
    (X : RSVPField R) (h : PositiveAccess X)
    (U : Region R) ((*$\varepsilon$*) : Real) : Prop :=
  (*$\forall$*) x (*$\in$*) U.carrier, |chi X h x - 1| < (*$\varepsilon$*)
\end{lstlisting}

\subsection{Reconstruction Sheaf Structure}

\begin{lstlisting}[caption={Lean-style reconstruction sheaf}]
structure ReconSheaf (R : ReconDomain) (X : RSVPField R) where
  sections : Region R -> Type
  restrict :
    (*$\forall$*) {U V : Region R},
      V.carrier (*$\subseteq$*) U.carrier ->
      sections U -> sections V
  identity_law     : Prop  -- restrict id = id
  composition_law  : Prop  -- restrict (V (*$\subseteq$*) U) after (W (*$\subseteq$*) V) =
                            --   restrict (W (*$\subseteq$*) U)
  gluing_law       : Prop  -- compatible sections glue
\end{lstlisting}

\subsection{Obstruction Formalization}

\begin{lstlisting}[caption={Compatibility and dispersion obstruction in Lean}]
def CompatObstruction {R : ReconDomain} {X : RSVPField R}
    (F : ReconSheaf R X)
    (U V : Region R)
    (sU : F.sections U) (sV : F.sections V) : Prop :=
  -- restrictions to overlap disagree
  (*$\neg$*)(F.restrict (overlap_le U V) sU =
     F.restrict (overlap_le' U V) sV)

structure Scale where
  ell : Real
  pos : ell > 0

-- Persistence functional placeholder
def Persistence {R : ReconDomain} {X : RSVPField R}
    (F : ReconSheaf R X) ((*$\ell$*) : Scale) (s : Type) : Real :=
  0 -- to be derived from field equations

def DispersionObstruction {R : ReconDomain} {X : RSVPField R}
    (F : ReconSheaf R X) ((*$\varepsilon$*) : Real) (s : Type) : Prop :=
  (*$\exists$*) (*$\ell$*) : Scale, Persistence F (*$\ell$*) s < (*$\varepsilon$*)
\end{lstlisting}

\subsection{Theorem Schemas in Lean Style}

\begin{lstlisting}[caption={Central theorem schemas as Lean proof obligations}]
-- Entropy convergence schema
theorem entropy_convergence
    {R : ReconDomain} (X : RSVPField R)
    (h : PositiveAccess X)
    (F : ReconSheaf R X)
    (U : Region R) ((*$\delta$*) (*$\varepsilon$*) : Real)
    (hFact : StronglyFactorizable X h U (*$\delta$*))
    (hBound : (* dispersion obstruction bounded *) True) :
    (* entropy channels converge to order (*$\varepsilon$*) *)
    True := by
  sorry  -- proof obligation: requires field equation dynamics

-- Entropy divergence schema
theorem entropy_divergence
    {R : ReconDomain} (X : RSVPField R)
    (h : PositiveAccess X)
    (F : ReconSheaf R X)
    (U : Region R) ((*$\varepsilon$*) : Real)
    (hTrans : NearTransition X h U (*$\varepsilon$*)) :
    (* entropy channels diverge under coarse-graining *)
    True := by
  sorry  -- proof obligation: requires beta-function derivation

-- Central conjecture as proof obligation
-- (RA proportional to c/H_0 in weak-gradient limit)
theorem RA_scaling_conjecture
    (field_eqs : True) :
    True := by
  sorry  -- principal open problem of the paper
\end{lstlisting}

The \texttt{sorry} markers are not weaknesses: they identify proof obligations. A
Lean formalization would force every informal claim in the paper to become either a
definition, theorem, conjecture, or unresolved proof target --- precisely the
discipline the present paper maintains in prose.

\subsection{Algorithmic Reconstruction Procedure}

A computational implementation of RSVP reconstruction obstruction analysis:

\begin{lstlisting}[caption={RSVP obstruction analysis algorithm}]
Algorithm: RSVP_Obstruction_Analysis

Input:
  X = (Phi, v, S)          -- coupled RSVP field
  R                         -- reconstruction domain
  U = {U_alpha}             -- initial admissibility cover at l_min
  l_min, l_max              -- scale range
  epsilon                   -- stability threshold
  delta                     -- factorization threshold

Output:
  R_A                       -- admissibility radius
  compat_obs_map            -- compatibility obstruction map
  disp_obs_map              -- dispersion obstruction map
  entropy_divergence_profile -- entropy channel divergence

Procedure:
  1. Initialize cover U at scale l = l_min.
  2. For each U_alpha:
       compute chi = S / Phi
       if chi < delta: classify as strongly factorizable
  3. For each overlap U_alpha cap U_beta:
       restrict local sections to overlap
       test compatibility agreement
       if disagreement: record compatibility obstruction
  4. Coarse-grain cover to next scale l.
  5. Update Phi_l, v_l, S_l.
  6. Recompute chi_l = S_l / Phi_l.
  7. Track persistence of previously admissible sections.
  8. If compatible sections fail persistence: record disp. obstruction.
  9. If chi_l -> 1 or persistence < epsilon:
       record R_A = current scale
  10. Compute entropy projections Pi_therm, Pi_info, Pi_grav, Pi_conf.
  11. Measure entropy-channel convergence or divergence.
  12. Repeat steps 4-11 until l = l_max.
  13. Return obstruction and transition profiles.
\end{lstlisting}

\subsection{Formalization Roadmap}

A rigorous implementation program would proceed in stages:
\begin{enumerate}[label=(\arabic*)]
  \item Define the typed syntax and semantics of RSVP reconstruction objects.
  \item Implement the reconstruction sheaf in Lean 4 or Agda.
  \item Formalize compatibility obstruction through \v{C}ech cocycle data.
  \item Define dispersion obstruction as scale-instability of persistence.
  \item Implement coarse-graining flow for $\chi_\ell$.
  \item Prove conditional theorems linking strong factorization to entropy convergence.
  \item Prove or refute the central conjecture: weak-gradient RSVP determines
    $\RA \propto c/H_0$.
  \item Build numerical simulations computing obstruction maps and entropy-channel
    divergence profiles from candidate field equations.
\end{enumerate}

The purpose of formalization is not merely aesthetic rigor. It is a guardrail against
the main failure mode of speculative gravity frameworks: sliding between definitions,
metaphors, conjectures, and derivations without marking the transition. A Lean
implementation would force every RSVP claim to declare its type.

\newpage
% ── App I ─────────────────────────────────────────────────────────────────────
\section{Toward a Variational Formulation}
\label{app:variational}

The candidate dynamical equations of Appendix~\ref{app:dynamics} appear
phenomenological. A variational formulation would dramatically increase mathematical
maturity and reveal whether the RSVP admissibility structure has a natural Lagrangian
origin. This appendix sketches the form such a formulation would take without
claiming to complete it.

\subsection{Candidate Action}

Introduce a candidate action functional over the reconstruction domain $\recon$:
\begin{equation}
  \label{eq:action}
  \mathcal{S}[X] \;=\; \int_{\recon}
    \mathcal{L}\!\left(\Phi,\,\mathbf{v},\,S,\,\nabla\Phi,\,\nabla\mathbf{v},\,\nabla S\right)
    \,d\mu,
\end{equation}
where $d\mu$ is a measure on $\recon$ derived from the admissibility geometry rather
than a fixed background volume form.

A minimal candidate Lagrangian density consistent with the dimensional analysis of
Section~\ref{sec:dimensional} and the qualitative dynamics of
Appendix~\ref{app:dynamics} is:
\begin{equation}
  \label{eq:lagrangian}
  \mathcal{L} \;=\;
    \frac{1}{2}(\nabla\Phi)^2
    \;-\; V(\chi)\,\Phi^2
    \;+\; \frac{1}{2}\,\mathbf{v}\cdot\mathbf{v}
    \;-\; \mathcal{W}(S,\Phi,\mathbf{v}),
\end{equation}
where $V(\chi) = V(S/\Phi)$ is a potential encoding the admissibility phase
structure and $\mathcal{W}$ is a coupling term generating the interaction between
dispersion growth and coherence transport.

The potential $V(\chi)$ must satisfy:
\begin{itemize}
  \item $V(0) = 0$: no potential energy at the Newtonian fixed point.
  \item $V'(\chi) > 0$: potential increases toward the transition, disfavoring
    high-dispersion configurations.
  \item A local minimum near $\chi = 0$ and a saddle or maximum near $\chi = 1$,
    consistent with the attractor structure of Appendix~\ref{app:attractor}.
\end{itemize}

\subsection{Admissibility Persistence as Modified Stationarity}

The central insight distinguishing the RSVP variational formulation from standard
Euler--Lagrange theory is that admissibility persistence is \emph{not} equivalent to
stationarity of $\mathcal{S}$.

Standard field theory requires:
\[
  \frac{\delta \mathcal{S}}{\delta X} \;=\; 0.
\]
RSVP admissibility persistence additionally requires that the stationary
configuration remain stable under renormalization of the reconstruction cover:
\[
  \frac{\delta \mathcal{S}}{\delta X} = 0
  \quad \text{and} \quad
  \mathcal{P}[\delta X / \delta \ell] > \varepsilon
  \quad \text{for all } \ell < \RA.
\]
The second condition filters out stationary configurations that are classically
extremal but fail reconstruction persistence under scale extension --- the
field-theoretic analog of dispersion obstruction.

This distinction reinforces the paper's central inversion: the relevant selection
criterion for physical configurations is not classical stationarity but admissibility
persistence under coarse-graining of the reconstruction cover.

\begin{obligation}
\label{obl:variational}
Derive the coupled field equations of Appendix~\ref{app:dynamics} as Euler--Lagrange
equations of a specific action of the form~\eqref{eq:action}, and show that the
admissibility persistence condition is compatible with, but strictly stronger than,
classical stationarity.
\end{obligation}

% ── App J ─────────────────────────────────────────────────────────────────────
\section{Effective Metric Emergence}
\label{app:metric}

Spacetime geometry is discussed philosophically in the main text but not yet
mathematically. This appendix sketches how an effective metric structure might emerge
from the admissibility geometry of the coupled field $X = (\Phi, \mathbf{v}, S)$
without presupposing a background metric. The effective metric construction developed
here is heuristic and should be interpreted as an emergent reconstruction observable
rather than a completed relativistic metric theory. Writing $g_{\mu\nu}^{\mathrm{eff}}$
immediately raises questions of covariance, constraint algebra, stress-energy
conservation, hyperbolicity, and equivalence principle; these are not resolved here,
and the appendix should be read as identifying the structural direction rather than
claiming completion.

\subsection{Effective Metric from Admissibility Structure}

In strongly factorizable regimes ($\chi \ll 1$), the continuation accessibility
field $\Phi$ and coherence-flow field $\mathbf{v}$ together define a preferred local
geometry. A schematic effective metric may be constructed as:
\begin{equation}
  \label{eq:effmetric}
  g_{\mu\nu}^{\mathrm{eff}} \;=\; G_{\mu\nu}(\Phi, \mathbf{v}, S),
\end{equation}
where $G_{\mu\nu}$ is a symmetric tensor built from the gradients of the RSVP fields.
The simplest candidate consistent with dimensional constraints is:
\[
  g_{\mu\nu}^{\mathrm{eff}} \;=\; f(\chi)\,\eta_{\mu\nu}
    \;+\; h(\chi)\,\frac{\partial_\mu\Phi\,\partial_\nu\Phi}{\Phi^2},
\]
where $\eta_{\mu\nu}$ is a flat reference metric, $f$ and $h$ are functions of
$\chi$ satisfying $f(0) = 1$ and $h(0) = 0$ (recovering flat geometry in the
Newtonian limit), and corrections from $h$ encode the admissibility curvature
generated by $\Phi$-gradients.

\subsection{Null Structure and Effective Geodesics}

Within this framework, null geodesics --- paths along which $g_{\mu\nu}^{\mathrm{eff}}
\,dx^\mu dx^\nu = 0$ --- would correspond to admissible photon continuation trajectories.
In the Newtonian regime ($f \approx 1$, $h \approx 0$), these reduce to ordinary
null geodesics of flat spacetime. Near the admissibility transition ($\chi \to 1$),
the $h$-term introduces curvature driven by accessibility gradients $\nabla\Phi$,
producing lensing effects without requiring a background metric curvature sourced by
matter.

This provides a schematic geometric basis for the lensing reinterpretation of
Section~\ref{sec:darkmatter}: lensing anomalies arise from admissibility curvature in
$g_{\mu\nu}^{\mathrm{eff}}$ rather than from hidden mass distributions.

\subsection{Relationship to General Relativity}

The effective metric~\eqref{eq:effmetric} is not postulated to satisfy the Einstein
field equations. Instead, the Einstein equations would appear as an emergent
approximation valid in regimes where the RSVP field dynamics produce a metric
satisfying the vacuum or sourced Einstein equations to leading order. This would
require showing that the Ricci tensor of $g_{\mu\nu}^{\mathrm{eff}}$ satisfies
$R_{\mu\nu} - \frac{1}{2}g_{\mu\nu}R \approx 8\pi G\, T_{\mu\nu}^{\mathrm{eff}}$
in appropriate regimes. This constitutes an open derivation target.

\begin{obligation}
\label{obl:metric}
Derive the explicit form of $G_{\mu\nu}(\Phi, \mathbf{v}, S)$ from the RSVP field
equations, compute the resulting Riemann curvature tensor, and determine in which
reconstruction regimes the effective geometry approximates solutions of the Einstein
equations.
\end{obligation}

% ── App K ─────────────────────────────────────────────────────────────────────
\section{Gauge Redundancy and Admissibility Equivalence}
\label{app:gauge}

Since RSVP increasingly resembles a geometric theory rather than a collection of
fields, the question of gauge redundancy arises naturally: do distinct field
configurations $X$ and $X'$ represent the same physical reconstruction geometry?

\subsection{Admissibility Equivalence}

Define an \emph{admissibility equivalence} relation on field configurations:
\[
  X \;\sim\; X'
\]
if and only if all entropy projection functors agree:
\[
  \Pi_i(X) \;\simeq\; \Pi_i(X') \quad \text{for all } i \in \{\mathrm{therm},
  \mathrm{info}, \mathrm{grav}, \mathrm{conf}\},
\]
and the obstruction classes of the associated reconstruction sheaves are identical:
\[
  \Hcech^k(\Ucover_X, \sheaf_X) \;\cong\; \Hcech^k(\Ucover_{X'}, \sheaf_{X'})
  \quad \text{for all } k.
\]
Two configurations that are admissibility-equivalent produce the same entropy-channel
structure, the same obstruction classes, and therefore the same physical predictions.
The physical state space of RSVP is the quotient:
\[
  \mathcal{M}_{\mathrm{phys}} \;=\; \{X\} \;/\; {\sim}.
\]

\subsection{Reconstruction-Preserving Transformations}

A \emph{reconstruction-preserving transformation} is a map $\phi: X \mapsto X'$
such that $X \sim X'$. These transformations form a group $\mathcal{G}_{\mathrm{adm}}$
--- the admissibility gauge group --- whose structure encodes the redundancy of the
field description.

In the Newtonian regime ($\chi \ll 1$), $\mathcal{G}_{\mathrm{adm}}$ likely reduces
to familiar diffeomorphism-like redundancy: smooth reparametrizations of
reconstruction neighborhoods that preserve admissibility classes. Near the
admissibility transition, the gauge group structure may change, reflecting the
breakdown of factorization.

\subsection{Relationship to Diffeomorphism Invariance}

General relativity's diffeomorphism invariance is the statement that the physical
content of a spacetime metric is invariant under smooth coordinate changes. In RSVP,
the analogous invariance is admissibility equivalence: physical content is invariant
under reconstruction-preserving transformations that preserve the sheaf structure and
entropy projections.

This suggests that RSVP's gauge group generalizes diffeomorphism invariance: in
strongly factorizable regimes, the two should approximately coincide, while near the
admissibility transition, the gauge group may become larger, encoding the additional
freedom in how cosmological boundary data is incorporated into local reconstruction.

\begin{obligation}
\label{obl:gauge}
Classify the admissibility gauge group $\mathcal{G}_{\mathrm{adm}}$ in both the
Newtonian regime ($\chi \ll 1$) and near the admissibility transition ($\chi \to 1$),
and determine whether RSVP possesses a well-posed constraint surface analogous to the
Gauss-law and diffeomorphism constraints of general relativity.
\end{obligation}

\newpage
% ══════════════════════════════════════════════════════════════════════════════
% Bibliography
% ══════════════════════════════════════════════════════════════════════════════

\begin{thebibliography}{99}

% ── Entropic / thermodynamic gravity ─────────────────────────────────────────
\bibitem{kastner2024}
R.~E. Kastner and A.~Schlatter,
``Gravity from Transactions: Fulfilling the Entropic Gravity Program,''
\textit{preprint} (2024).

\bibitem{verlinde2011}
E.~Verlinde,
``On the Origin of Gravity and the Laws of Newton,''
\textit{J.\ High Energy Phys.}\ \textbf{2011}, 29 (2011).

\bibitem{verlinde2017}
E.~Verlinde,
``Emergent Gravity and the Dark Universe,''
\textit{SciPost Phys.}\ \textbf{2}, 016 (2017).

\bibitem{jacobson1995}
T.~Jacobson,
``Thermodynamics of Spacetime: The Einstein Equation of State,''
\textit{Phys.\ Rev.\ Lett.}\ \textbf{75}, 1260 (1995).

\bibitem{jacobson2016}
T.~Jacobson,
``Entanglement Equilibrium and the Einstein Equation,''
\textit{Phys.\ Rev.\ Lett.}\ \textbf{116}, 201101 (2016).

\bibitem{padmanabhan2010}
T.~Padmanabhan,
``Thermodynamical Aspects of Gravity: New Insights,''
\textit{Rep.\ Prog.\ Phys.}\ \textbf{73}, 046901 (2010).

% ── MOND and modified gravity ─────────────────────────────────────────────────
\bibitem{milgrom1983}
M.~Milgrom,
``A Modification of the Newtonian Dynamics as a Possible Alternative to the Hidden
Mass Hypothesis,''
\textit{Astrophys.\ J.}\ \textbf{270}, 365 (1983).

\bibitem{bekenstein2004}
J.~D. Bekenstein,
``Relativistic Gravitation Theory for the Modified Newtonian Dynamics Paradigm,''
\textit{Phys.\ Rev.\ D}\ \textbf{70}, 083509 (2004).

\bibitem{famaey2012}
B.~Famaey and S.~McGaugh,
``Modified Newtonian Dynamics (MOND): Observational Phenomenology and Relativistic
Extensions,''
\textit{Living Rev.\ Relativ.}\ \textbf{15}, 10 (2012).

\bibitem{mcgaugh2011}
S.~S. McGaugh,
``A Novel Test of the Modified Newtonian Dynamics with Gas Rich Galaxies,''
\textit{Phys.\ Rev.\ Lett.}\ \textbf{106}, 121303 (2011).

% ── Holographic / entanglement spacetime ─────────────────────────────────────
\bibitem{vanraamsdonk2010}
M.~Van Raamsdonk,
``Building Up Spacetime with Quantum Entanglement,''
\textit{Gen.\ Relativ.\ Gravit.}\ \textbf{42}, 2323 (2010).

\bibitem{swingle2012}
B.~Swingle,
``Entanglement Renormalization and Holography,''
\textit{Phys.\ Rev.\ D}\ \textbf{86}, 065007 (2012).

% ── Sheaf theory and algebraic topology ──────────────────────────────────────
\bibitem{bredon1997}
G.~E. Bredon,
\textit{Sheaf Theory}, 2nd ed.
Springer, New York, 1997.

\bibitem{kashiwara2005}
M.~Kashiwara and P.~Schapira,
\textit{Categories and Sheaves}.
Springer, Berlin, 2005.

\bibitem{maclane1992}
S.~Mac Lane and I.~Moerdijk,
\textit{Sheaves in Geometry and Logic}.
Springer, New York, 1992.

\bibitem{bott1982}
R.~Bott and L.~W. Tu,
\textit{Differential Forms in Algebraic Topology}.
Springer, New York, 1982.

\bibitem{hatcher2002}
A.~Hatcher,
\textit{Algebraic Topology}.
Cambridge University Press, Cambridge, 2002.

% ── Renormalization and phase transitions ────────────────────────────────────
\bibitem{wilson1974}
K.~G. Wilson and J.~Kogut,
``The Renormalization Group and the $\epsilon$ Expansion,''
\textit{Phys.\ Rep.}\ \textbf{12}, 75 (1974).

\bibitem{goldenfeld1992}
N.~Goldenfeld,
\textit{Lectures on Phase Transitions and the Renormalization Group}.
Addison-Wesley, Reading, MA, 1992.

% ── Entropy and information geometry ─────────────────────────────────────────
\bibitem{jaynes1957}
E.~T. Jaynes,
``Information Theory and Statistical Mechanics,''
\textit{Phys.\ Rev.}\ \textbf{106}, 620 (1957).

\bibitem{amari2016}
S.~Amari,
\textit{Information Geometry and Its Applications}.
Springer, Tokyo, 2016.

\end{thebibliography}

\end{document}
